% !TEX program = xelatex
%% =====================================================================
%%  LTspice Laboratory Manual  -  Electronic Devices and Circuits (CO2037)
%%  Dự án 3: Xây dựng tài liệu thực hành mới với LTspice
%%  Biên dịch: xelatex LTspice_Manual.tex   (chạy 2 lần)   hoặc   latexmk LTspice_Manual.tex
%% =====================================================================
\documentclass[12pt,twoside]{Book}

%% ------------- Thông tin nhóm biên soạn (sửa tại đây) -------------
\newcommand{\GroupNo}{9}
\newcommand{\MemberA}{Phạm Đỗ Hồng Phúc}   \newcommand{\IdA}{2510359}
\newcommand{\MemberB}{Lê Lương Vỹ}          \newcommand{\IdB}{2214056}
\newcommand{\MemberC}{Nguyễn Bảo Duy Thành} \newcommand{\IdC}{2510372}
%% -------------------------------------------------------------------

%% =====================================================================
%%  manual_common.tex - cấu hình chung của tài liệu thực hành LTspice
%%  (bố cục theo Book.cls và các logbook; tiếng Anh, font Utopia thu 0.92)
%%  Biên dịch: xelatex LTspice_Manual.tex (2 lần)
%% =====================================================================
\ifdefined\directlua\usepackage{fontspec}\else\ifdefined\XeTeXversion\usepackage{fontspec}\fi\fi
\usepackage[a4paper,tmargin=2cm,bmargin=2.5cm,lmargin=3.5cm,rmargin=2.0cm]{geometry}
\usepackage{color,calc,graphicx}
\usepackage{amsmath,amssymb,amsfonts,enumitem}
\usepackage{gensymb}
\providecommand{\micro}{\ifmmode\text{\textmu}\else\textmu\fi}
\usepackage{float,titlesec,fancyhdr,chngcntr,eso-pic,afterpage,array,booktabs}
\usepackage[table]{xcolor}
\usepackage[most]{tcolorbox}
\usepackage[T1]{fontenc}
\makeatletter
\def\pus@decl#1#2#3#4#5{\DeclareFontShape{#1}{pus}{#2}{#3}{<->s*[0.92]#4}{}}
\DeclareFontFamily{T1}{pus}{}
\pus@decl{T1}{m}{n}{putr8t}{}\pus@decl{T1}{m}{sl}{putro8t}{}\pus@decl{T1}{m}{it}{putri8t}{}
\pus@decl{T1}{m}{sc}{putrc8t}{}
\pus@decl{T1}{b}{n}{putb8t}{}\pus@decl{T1}{b}{sl}{putbo8t}{}\pus@decl{T1}{b}{it}{putbi8t}{}
\DeclareFontShape{T1}{pus}{bx}{n}{<->ssub * pus/b/n}{}
\DeclareFontShape{T1}{pus}{bx}{sl}{<->ssub * pus/b/sl}{}
\DeclareFontShape{T1}{pus}{bx}{it}{<->ssub * pus/b/it}{}
\DeclareFontFamily{TS1}{pus}{}
\pus@decl{TS1}{m}{n}{putr8c}{}\pus@decl{TS1}{m}{sl}{putro8c}{}\pus@decl{TS1}{m}{it}{putri8c}{}
\pus@decl{TS1}{b}{n}{putb8c}{}\pus@decl{TS1}{b}{sl}{putbo8c}{}\pus@decl{TS1}{b}{it}{putbi8c}{}
\DeclareFontShape{TS1}{pus}{bx}{n}{<->ssub * pus/b/n}{}
\DeclareFontShape{TS1}{pus}{bx}{sl}{<->ssub * pus/b/sl}{}
\DeclareFontShape{TS1}{pus}{bx}{it}{<->ssub * pus/b/it}{}
\DeclareFontFamily{OT1}{pus}{}
\pus@decl{OT1}{m}{n}{putr7t}{}\pus@decl{OT1}{m}{sl}{putro7t}{}\pus@decl{OT1}{m}{it}{putri7t}{}
\pus@decl{OT1}{b}{n}{putb7t}{}\pus@decl{OT1}{b}{sl}{putbo7t}{}\pus@decl{OT1}{b}{it}{putbi7t}{}
\DeclareFontShape{OT1}{pus}{bx}{n}{<->ssub * pus/b/n}{}
\DeclareFontShape{OT1}{pus}{bx}{sl}{<->ssub * pus/b/sl}{}
\DeclareFontShape{OT1}{pus}{bx}{it}{<->ssub * pus/b/it}{}
\makeatother
\renewcommand{\ttdefault}{pcr}
\renewcommand{\rmdefault}{pus}\renewcommand{\familydefault}{\rmdefault}
\usepackage[italic]{mathastext}
\ifdefined\directlua\newfontfamily\vnfont{Latin Modern Roman}\else\ifdefined\XeTeXversion\newfontfamily\vnfont{Latin Modern Roman}\fi\fi
\usepackage[colorlinks=true,linkcolor=black,urlcolor=blue]{hyperref}

\definecolor{accent}{HTML}{D98200}
\definecolor{lightorange}{HTML}{FFF3DC}
\definecolor{navy}{HTML}{1F2F5E}
\definecolor{lightnavy}{HTML}{EAEFF8}

\pagestyle{fancy}
\lhead[]{}\chead[]{}\rhead[]{}
\lfoot[\sl Page \thepage]{\it LTspice Laboratory Manual \hfil}
\cfoot[]{}
\rfoot[\it HCMUT - Computer Engineering]{\sl Page \thepage}
\renewcommand{\headrulewidth}{0pt}\renewcommand{\footrulewidth}{1pt}
\setlength{\parindent}{0pt}
\newcommand\blankpage{\null\thispagestyle{empty}\addtocounter{page}{0}\newpage}

\renewcommand{\chaptername}{Lab}
\titleformat{\chapter}[display]{\bfseries\huge}{\filright\MakeUppercase{\chaptertitlename} \huge\thechapter}{21.8pt}{\titlerule\vspace{11.05pt}\filleft}[\vspace{10.85pt}\titlerule]
\titlespacing*{\chapter}{0pt}{50.1pt}{40pt}
\counterwithin*{section}{chapter}\counterwithin*{subsection}{section}\counterwithin*{subsubsection}{subsection}
\renewcommand{\chap}[1]{\chapter{#1}\thispagestyle{empty}%
  \AddToShipoutPictureBG*{\includegraphics[width=\paperwidth,height=\paperheight]{images/template/chapter.png}}\newpage}

% ---------- Dotted answer lines (calculation parts are left blank for students) ----------
\newcount\dotcnt
\newcommand{\dotsx}[3]{\par\vspace{#1}\begingroup\parskip=0pt \baselineskip=#2 \lineskip=0pt \lineskiplimit=0pt
  \dotcnt=#3\relax \loop\ifnum\dotcnt>0 \noindent\dotfill\par\advance\dotcnt by -1 \repeat\endgroup}
\newcommand{\dotlines}[1]{\dotsx{0pt}{26.5pt}{#1}}
\newcommand{\notesblock}{\textbf{\textit{Notes:}}\\
  \textit{Explanations, formulas, and equations are expected rather than only results.}\par}
\newcommand{\theory}[1]{\subsubsection{Theory calculation}\vspace{-3pt}\notesblock\dotsx{2.4pt}{26.5pt}{#1}}

% ---------- Placeholders / figures (LTspice screenshots go to images/ or images/LabN/) ----------
\graphicspath{{images/}{images/template/}{images/Lab1/}{images/Lab2/}{images/Lab3/}{images/Lab4/}{images/Lab5/}{images/Lab6/}}

\newcommand{\placeholderbox}[2]{% file, width
  \begin{tcolorbox}[width=#2,colback=lightorange,colframe=accent,boxrule=1pt,arc=2pt,halign=center,valign=center,height=3.2cm,center]
  \color{accent!85!black}\bfseries\itshape [Insert LTspice screenshot]\\[1pt]
  {\small\normalfont\upshape(save in \texttt{images/LabN/\detokenize{#1}} or \texttt{images/\detokenize{#1}})}\end{tcolorbox}}

\newcommand{\labfig}[4]{% file, width, caption, label
  \begin{figure}[H]\centering
  \IfFileExists{images/#1}{\includegraphics[width=#2]{images/#1}}{%
    \IfFileExists{images/Lab1/#1}{\includegraphics[width=#2]{images/Lab1/#1}}{%
    \IfFileExists{images/Lab2/#1}{\includegraphics[width=#2]{images/Lab2/#1}}{%
    \IfFileExists{images/Lab3/#1}{\includegraphics[width=#2]{images/Lab3/#1}}{%
    \IfFileExists{images/Lab4/#1}{\includegraphics[width=#2]{images/Lab4/#1}}{%
    \IfFileExists{images/Lab5/#1}{\includegraphics[width=#2]{images/Lab5/#1}}{%
    \IfFileExists{images/Lab6/#1}{\includegraphics[width=#2]{images/Lab6/#1}}{%
    \placeholderbox{#1}{#2}}}}}}}%
  \caption{#3}\label{#4}\end{figure}}

% ---------- Comparison table and simulation block ----------
\newcommand{\cmprow}[1]{#1 & & & \\ \hline}
\newcommand{\comparetable}[1]{\subsubsection{Comparison}
  \begin{center}\renewcommand{\arraystretch}{1.9}
  \begin{tabular}{|p{4.6cm}|c|c|c|}\hline
  \textbf{Quantity} & \textbf{Theory} & \textbf{LTspice} & \textbf{Error (\%)}\\ \hline #1
  \end{tabular}\end{center}}
\newcommand{\commentlines}[1]{\textbf{\textit{Comments:}} \dotfill\dotsx{0pt}{26.5pt}{#1}}
\newcommand{\simulation}[1]{\subsubsection{LTspice simulation}\begin{itemize}[itemsep=3pt,topsep=3pt]#1\end{itemize}}
\newcommand{\cmd}[1]{\texttt{#1}}
\newcommand{\status}[1]{{\small\noindent\textbf{\color{navy}Status:} \textit{#1}\par}\vspace{4pt}}
\emergencystretch=5em
\tolerance=2000
\newcommand{\file}[1]{\texttt{#1}}


\begin{document}

%% ----------------------------- Title page -----------------------------
\thispagestyle{empty}
\begin{center}
\vspace*{1.2cm}
{\bfseries\large HO CHI MINH CITY UNIVERSITY OF TECHNOLOGY\\
FACULTY OF COMPUTER SCIENCE AND ENGINEERING\par}
\vspace{1.6cm}
{\color{navy}\bfseries Electronic Devices and Circuits (CO2037)\par}
\vspace{1.2cm}
{\bfseries\Huge LTspice Laboratory Manual\par}
\vspace{6pt}
{\itshape\large Six laboratories on circuit simulation with LTspice\par}
\vspace{1.4cm}
{\large Student edition -- Laboratories 1 to 6\par}
\vspace{2.2cm}
\renewcommand{\arraystretch}{1.5}
\begin{tabular}{|p{5.8cm}|p{5.8cm}|}\hline
\textbf{Prepared by (Group \GroupNo)} & \textbf{Student ID}\\ \hline
{\vnfont\MemberA} & \IdA\\ \hline
{\vnfont\MemberB} & \IdB\\ \hline
{\vnfont\MemberC} & \IdC\\ \hline
\end{tabular}
\vfill
Instructor: Dr. Le Trong Nhan
\end{center}
\clearpage

%% ----------------------------- Preface -----------------------------
\section*{Preface}
This manual replaces the former PSpice for TI laboratories with LTspice, a free circuit simulator from Analog Devices. It contains six laboratories; each laboratory has ten main exercises and a set of advanced exercises.

\textbf{What you need.} The offline LTspice installation provided for the course (including the component libraries). All schematics in this manual open and run on that installation without downloading anything else.

\textbf{How to work.}
\begin{enumerate}
  \item Read the introduction and the LTspice guide of the laboratory.
  \item For each exercise, complete the \emph{Theory calculation} by hand: the calculation space is intentionally blank, and the reference solutions are given only to the instructor.
  \item Build the circuit in LTspice, run the analysis indicated in the \emph{LTspice simulation} part and save the schematic as \file{LabN/ExNN.asc} (advanced exercises: \file{LabN/AdvNN.asc}).
  \item Fill in the \emph{Comparison} table with your theoretical and simulated values and comment on any difference.
\end{enumerate}

\textbf{Conventions.} Component names follow the figures (R1, V1, \dots). Values use SPICE prefixes (\texttt{k}, \texttt{Meg}, \texttt{m}, \texttt{u}); remember that \texttt{M} also means milli in SPICE. Voltage and current signs follow the reference directions shown in each figure; LTspice reports the current of a voltage source as positive when it enters the positive terminal.
\clearpage

\tableofcontents
\clearpage

\chap{DC Circuit Analysis -- First Project on LTspice}

\section{Introduction}
This first laboratory introduces LTspice, the circuit simulator used throughout this manual. You will build resistive circuits, run a DC operating-point analysis and compare the simulated values with your own hand calculations.

\textbf{After this laboratory you will be able to:}
\begin{itemize}
  \item create, wire and save a schematic in LTspice;
  \item set component values with SPICE unit prefixes and place ideal sources and dependent sources;
  \item run a DC operating-point analysis (\cmd{.op}) and read node voltages and branch currents;
  \item verify Kirchhoff's laws, equivalent resistance, circuit transformations and power balance by simulation;
  \item compare theory with simulation and explain the differences;
  \item use parameters (\cmd{.param}), parametric steps (\cmd{.step}) and measurements (\cmd{.meas}) to automate the comparison.
\end{itemize}

\textbf{How the exercises are organised.} Every exercise has the same structure: the statement and the circuit figure; a \emph{Theory calculation} part that you complete by hand; an \emph{LTspice simulation} part that tells you what to build and which analysis to run; and a \emph{Comparison} table where you put your results side by side. The calculation space is intentionally left blank. Save every schematic as \file{Lab1/ExNN.asc}, where \texttt{NN} is the exercise number.

\section{LTspice guide}
\subsection{Create a schematic}
Start LTspice and choose \textit{File $\rightarrow$ New Schematic} (Figure~\ref{fig:lab1_guide01}). Save it immediately (\textit{File $\rightarrow$ Save As}) as \file{Lab1/Ex01.asc}.
\labfig{lab1_guide01.png}{10cm}{Create a new schematic in LTspice}{fig:lab1_guide01}

\subsection{Place and wire components}
\begin{itemize}
  \item Press \texttt{R} to place a resistor, \texttt{G} to place the ground (node \texttt{0}). Press \texttt{F2} to open the component list: search \cmd{voltage} for a DC voltage source, \cmd{current} for a current source, and \cmd{e}, \cmd{f}, \cmd{g}, \cmd{h} for the dependent sources.
  \item Rotate a component with \texttt{Ctrl+R} and mirror it with \texttt{Ctrl+E} before clicking to place it.
  \item Press \texttt{F3} (or the pencil icon) to draw wires between component pins. Press \texttt{F5} to delete and \texttt{F7} to move an object.
  \item Press \texttt{F4} to give a net a name (\textit{Label Net}), for example \texttt{A}, \texttt{B}.
  \item \textbf{Every circuit needs a ground} (node \texttt{0}); otherwise the simulation cannot run.
\end{itemize}
\labfig{lab1_guide02.png}{12cm}{Placing components and wires}{fig:lab1_guide02}

\subsection{Set the component values}
Right-click a component to open its dialog and type the value; for a voltage source enter the DC value. Use SPICE prefixes: \texttt{f}, \texttt{p}, \texttt{n}, \texttt{u} (micro), \texttt{m} (milli), \texttt{k}, \texttt{Meg}, \texttt{G}. \textbf{Warning:} in SPICE \texttt{m} and \texttt{M} both mean milli; mega must be written \texttt{Meg}.
\labfig{lab1_guide03.png}{9cm}{Editing the value of a component}{fig:lab1_guide03}

\subsection{Add the simulation command and run}
Choose \textit{Simulate $\rightarrow$ Edit Simulation Cmd}, select the \textit{DC op pnt} tab and click OK; the directive \cmd{.op} is placed on the schematic (you can also place it as text with the \texttt{S} key). Then choose \textit{Simulate $\rightarrow$ Run}.
\labfig{lab1_guide04.png}{10cm}{Adding the \texttt{.op} directive}{fig:lab1_guide04}

\subsection{Read the results}
After a \cmd{.op} run, a window lists the node voltages and the branch currents (Figure~\ref{fig:lab1_guide05}). If you use \cmd{.meas}, its results are written in the SPICE Error Log (\textit{View $\rightarrow$ SPICE Error Log}). Note the sign conventions: the current of a voltage source is positive when it flows into its positive terminal, so a source that delivers power shows a negative current.
\labfig{lab1_guide05.png}{10cm}{Operating-point results}{fig:lab1_guide05}

\section{Exercises}

\subsection{Exercise 1: Kirchhoff's voltage law with a dependent source}
Consider the circuit in Figure~\ref{fig:lab1_ex01}. It contains a 12\,V source, a 4\,V source, a 4\,$\Omega$ resistor, a 6\,$\Omega$ resistor and a voltage-controlled voltage source of value $2v_0$ (positive terminal on the left), where $v_0$ is the voltage across the 6\,$\Omega$ resistor (positive terminal on the left). Calculate the voltage $v_0$ and the current $i$ using Kirchhoff's voltage law (KVL). Then simulate the circuit to check your results.

\labfig{lab1_ex01.png}{7cm}{Exercise 1: circuit with a dependent source}{fig:lab1_ex01}

\theory{8}

\simulation{
\item Draw the circuit of Figure~\ref{fig:lab1_ex01} in LTspice (components \cmd{res}/\cmd{voltage}/\cmd{current}, ground \cmd{0}) and save it as \file{Lab1/Ex01.asc}.
\item Use the voltage-controlled voltage source \cmd{e} (press \texttt{F2} and search for \cmd{e}). Set its gain to 2 and connect its two control pins across the 6\,$\Omega$ resistor, with the polarity of $v_0$.
\item Add the \cmd{.op} analysis (\textit{Simulate $\rightarrow$ Edit Simulation Cmd $\rightarrow$ DC op pnt}), run the simulation and read the node voltages and branch currents in the operating-point window.
\item Be careful with signs: LTspice reports the current of a voltage source as positive when it flows \emph{into} its positive terminal, and the current of a resistor as positive when it flows from its first pin to its second pin.
}

\comparetable{
\cmprow{$v_0$ (V)}
\cmprow{$i$ (A)}
}

\commentlines{2}

\subsection{Exercise 2: Equivalent resistance by series and parallel reductions}
Consider the circuit in Figure~\ref{fig:lab1_ex02}. Rearrange the circuit to clarify its series and/or parallel topology, and draw the rearranged circuit. Then use series/parallel reductions to find the equivalent resistance seen by the source between terminals A and F. Finally, simulate the circuit to check that the current through the whole circuit is correctly calculated.

\labfig{lab1_ex02.png}{10cm}{Exercise 2: find the equivalent resistance between terminals A and F}{fig:lab1_ex02}

\theory{8}

\simulation{
\item Draw the circuit of Figure~\ref{fig:lab1_ex02} in LTspice (components \cmd{res}/\cmd{voltage}/\cmd{current}, ground \cmd{0}) and save it as \file{Lab1/Ex02.asc}.
\item Add the \cmd{.op} analysis (\textit{Simulate $\rightarrow$ Edit Simulation Cmd $\rightarrow$ DC op pnt}), run the simulation and read the node voltages and branch currents in the operating-point window.
\item Label the nodes A to F with \textit{Label Net} (\texttt{F4}), so that the results are listed with these names.
\item The equivalent resistance is $R_{AF}=V_1/I$, where $I$ is the magnitude of the current delivered by the source.
\item Be careful with signs: LTspice reports the current of a voltage source as positive when it flows \emph{into} its positive terminal, and the current of a resistor as positive when it flows from its first pin to its second pin.
}

\comparetable{
\cmprow{$R_{AF}$ ($\Omega$)}
\cmprow{$I$ (A)}
}

\commentlines{2}

\subsection{Exercise 3: Equivalent resistance and node voltages}
Consider the circuit in Figure~\ref{fig:lab1_ex03}. Rearrange the circuit to clarify its series and/or parallel topology. Next, find the equivalent resistance between terminals A and F and the voltages at A, B, C, D and E. Finally, simulate the circuit to check your calculation.

\labfig{lab1_ex03.png}{11cm}{Exercise 3: equivalent resistance and node voltages}{fig:lab1_ex03}

\theory{10}

\simulation{
\item Draw the circuit of Figure~\ref{fig:lab1_ex03} in LTspice (components \cmd{res}/\cmd{voltage}/\cmd{current}, ground \cmd{0}) and save it as \file{Lab1/Ex03.asc}.
\item Add the \cmd{.op} analysis (\textit{Simulate $\rightarrow$ Edit Simulation Cmd $\rightarrow$ DC op pnt}), run the simulation and read the node voltages and branch currents in the operating-point window.
\item Label the nodes A to F (\texttt{F4}); node F is the ground. The source is a current source: place it with \cmd{current} and check the arrow direction.
\item Be careful with signs: LTspice reports the current of a voltage source as positive when it flows \emph{into} its positive terminal, and the current of a resistor as positive when it flows from its first pin to its second pin.
}

\comparetable{
\cmprow{$R_{AF}$ ($\Omega$)}
\cmprow{$V_A$ (V)}
\cmprow{$V_B$ (V)}
\cmprow{$V_C$ (V)}
\cmprow{$V_D$ (V)}
\cmprow{$V_E$ (V)}
}

\commentlines{2}

\subsection{Exercise 4: Currents and voltages in a resistive network}
Given the circuit in Figure~\ref{fig:lab1_ex04}, find $I_1$, $I_2$, $I_3$, $V_a$ and $V_b$. Present your calculation steps and check them by simulation.

\labfig{lab1_ex04.png}{8cm}{Exercise 4: find $I_1$, $I_2$, $I_3$, $V_a$ and $V_b$}{fig:lab1_ex04}

\theory{8}

\simulation{
\item Draw the circuit of Figure~\ref{fig:lab1_ex04} in LTspice (components \cmd{res}/\cmd{voltage}/\cmd{current}, ground \cmd{0}) and save it as \file{Lab1/Ex04.asc}.
\item Add the \cmd{.op} analysis (\textit{Simulate $\rightarrow$ Edit Simulation Cmd $\rightarrow$ DC op pnt}), run the simulation and read the node voltages and branch currents in the operating-point window.
\item Compare the simulated branch currents with $I_1$, $I_2$, $I_3$ (check their reference directions).
\item Be careful with signs: LTspice reports the current of a voltage source as positive when it flows \emph{into} its positive terminal, and the current of a resistor as positive when it flows from its first pin to its second pin.
}

\comparetable{
\cmprow{$I_1$ (mA)}
\cmprow{$I_2$ (mA)}
\cmprow{$I_3$ (mA)}
\cmprow{$V_a$ (V)}
\cmprow{$V_b$ (V)}
}

\commentlines{2}

\subsection{Exercise 5: Selecting a standard resistor}
Given the network in Figure~\ref{fig:lab1_ex05}: a 24\,V source in series with an unknown resistor $R$ and a 3.9\,k$\Omega$ resistor, across which the voltage must be $V_1=5$\,V. Select resistor $R$ from the standard resistor list (table at the end of this lab) and complete the following requirements.

\labfig{lab1_ex05.png}{6cm}{Exercise 5: choose $R$ from the standard resistor list}{fig:lab1_ex05}

\textit{\textbf{Notes:}}\\
\textit{Explanations, formulas, and equations are expected rather than only results.}
\begin{enumerate}[label=\alph*.]
  \item Find the required value of the resistor $R$.\\ $R=$\dotfill
  \item Use the standard values of the table at the end of this lab to select a resistor with 10\,\% tolerance for $R$. $R$ may be a single resistor or a combination of resistors, as long as every resistor has a standard value and is available on the market.\\ The selected resistor:\dotfill
  \item Using the resistor selected in (b), determine the voltage across the 3.9\,k$\Omega$ resistor.\\ $V_1=$\dotfill
  \item Calculate the percent error in $V_1$ when the standard resistor of (b) is used.\\ \dotfill
  \item Determine the power rating for this standard component (choose among standard ratings such as 1/8, 1/4, 1/2 and 1\,W).\\ $P_R=$\dotfill
\end{enumerate}

\simulation{
\item Draw the circuit of Figure~\ref{fig:lab1_ex05} in LTspice (components \cmd{res}/\cmd{voltage}/\cmd{current}, ground \cmd{0}) and save it as \file{Lab1/Ex05.asc}.
\item Define the unknown resistor as a parameter: place a \cmd{.param} directive (\texttt{S}) such as \cmd{.param R=<value>} and type \cmd{\{R\}} in the value field of resistor R.
\item Add the \cmd{.op} analysis (\textit{Simulate $\rightarrow$ Edit Simulation Cmd $\rightarrow$ DC op pnt}), run the simulation and read the node voltages and branch currents in the operating-point window.
\item Use \cmd{.step param R list ...} with the selected value and its two neighbours in the standard list to see how $V_1$ changes (each step is one simulation run).
\item Be careful with signs: LTspice reports the current of a voltage source as positive when it flows \emph{into} its positive terminal, and the current of a resistor as positive when it flows from its first pin to its second pin.
}

\comparetable{
\cmprow{$R$ required ($\Omega$)}
\cmprow{$R$ selected ($\Omega$)}
\cmprow{$V_1$ (V)}
\cmprow{Error in $V_1$ (\%)}
\cmprow{$P_R$ (mW)}
}

\commentlines{2}

\subsection{Exercise 6: Circuit transformation (Y--$\Delta$)}
Consider the circuit in Figure~\ref{fig:lab1_ex06}. Apply what you have learned to transform it into another form in which the total equivalent resistance $R_{ab}$ is found more easily, and draw the transformed circuit. Next, find the current $i$ through the circuit and simulate it to check your results.

\labfig{lab1_ex06.png}{9cm}{Exercise 6: find $R_{ab}$ and the current $i$}{fig:lab1_ex06}

\theory{8}

\simulation{
\item Draw the circuit of Figure~\ref{fig:lab1_ex06} in LTspice (components \cmd{res}/\cmd{voltage}/\cmd{current}, ground \cmd{0}) and save it as \file{Lab1/Ex06.asc}.
\item Add the \cmd{.op} analysis (\textit{Simulate $\rightarrow$ Edit Simulation Cmd $\rightarrow$ DC op pnt}), run the simulation and read the node voltages and branch currents in the operating-point window.
\item $R_{ab}=V_1/i$; read $i$ from the current of the source (mind the sign).
\item Be careful with signs: LTspice reports the current of a voltage source as positive when it flows \emph{into} its positive terminal, and the current of a resistor as positive when it flows from its first pin to its second pin.
}

\comparetable{
\cmprow{$R_{ab}$ ($\Omega$)}
\cmprow{$i$ (A)}
}

\commentlines{2}

\subsection{Exercise 7: Circuit transformation and total current}
Consider the circuit in Figure~\ref{fig:lab1_ex07}. Transform it into another form in which the total equivalent resistance is found more easily, and draw the transformed circuit. Next, find the current $I_S$ through the circuit and simulate it to check your results.

\labfig{lab1_ex07.png}{6cm}{Exercise 7: find the equivalent resistance and $I_S$}{fig:lab1_ex07}

\theory{8}

\simulation{
\item Draw the circuit of Figure~\ref{fig:lab1_ex07} in LTspice (components \cmd{res}/\cmd{voltage}/\cmd{current}, ground \cmd{0}) and save it as \file{Lab1/Ex07.asc}.
\item Add the \cmd{.op} analysis (\textit{Simulate $\rightarrow$ Edit Simulation Cmd $\rightarrow$ DC op pnt}), run the simulation and read the node voltages and branch currents in the operating-point window.
\item Enter resistor values with the SPICE prefix \texttt{k} (for example \texttt{12k}); never use \texttt{M} for mega, because in SPICE \texttt{m} and \texttt{M} both mean milli (mega is \texttt{Meg}).
\item Be careful with signs: LTspice reports the current of a voltage source as positive when it flows \emph{into} its positive terminal, and the current of a resistor as positive when it flows from its first pin to its second pin.
}

\comparetable{
\cmprow{$R_{eq}$ (k$\Omega$)}
\cmprow{$I_S$ (mA)}
}

\commentlines{2}

\subsection{Exercise 8: Active and passive elements (1)}
In the circuit of Figure~\ref{fig:lab1_ex08}, the elements $p_2$, $p_3$ and $p_4$ are unknown elements; the voltage across each element and the current through it are given in the figure, together with a 30\,V source $p_1$ and a current-controlled current source $p_5$ of value $0.4I$. First, use the passive sign convention to identify whether each of $p_2$, $p_3$, $p_4$ \emph{absorbs} or \emph{supplies} power. Represent an absorbing element by a pure resistor with a proper value, and a supplying element by an ideal DC voltage source with the corresponding value. Next, redraw the circuit and calculate the power absorbed by each element (verify that the total power is zero). Then perform a simulation with the elements determined in the previous step.

\labfig{lab1_ex08.png}{10cm}{Exercise 8: determine the unknown elements and the absorbed power of each element}{fig:lab1_ex08}

\theory{8}

\simulation{
\item Draw the original circuit with $p_2$, $p_3$, $p_4$ as labelled rectangles (\textit{Draw $\rightarrow$ Rectangle}) and add the given voltages and currents as text. This drawing is only used for the statement; it is not simulated.
\item Redraw the circuit with the resistors/voltage sources you determined and save it as \file{Lab1/Ex08.asc}.
\item The current-controlled current source is \cmd{f}: its controlling current is the current of a voltage source (name it in the value field of \cmd{f}, together with the gain).
\item Add the \cmd{.op} analysis (\textit{Simulate $\rightarrow$ Edit Simulation Cmd $\rightarrow$ DC op pnt}), run the simulation and read the node voltages and branch currents in the operating-point window.
\item Compute the power of each element from the simulated voltages and currents.
\item Be careful with signs: LTspice reports the current of a voltage source as positive when it flows \emph{into} its positive terminal, and the current of a resistor as positive when it flows from its first pin to its second pin.
}

\comparetable{
\cmprow{$p_2$: type and power (W)}
\cmprow{$p_3$: type and power (W)}
\cmprow{$p_4$: type and power (W)}
\cmprow{$p_1$ (W)}
\cmprow{$p_5$ (W)}
}

\commentlines{2}

\subsection{Exercise 9: Active and passive elements (2)}
Consider the circuit in Figure~\ref{fig:lab1_ex09}, where a 10\,V source supplies a 12\,$\Omega$ resistor, the elements $p_1$ and $p_2$ are unknown, and a current-controlled voltage source of value $3i_x$ closes the loop; node D is the ground. Find the voltage $v$ and the current $i_x$. According to the result, determine whether the elements whose absorbed powers are $p_1$ and $p_2$ are active or passive (calculations are required); use the passive sign convention. Represent an element that consumes power by a pure resistor with an appropriate value, and a power-supplying element by a corresponding ideal DC voltage source. Perform a simulation to check how the circuit works.

\labfig{lab1_ex09.png}{8cm}{Exercise 9: find the unknown elements and variables}{fig:lab1_ex09}

\theory{10}

\simulation{
\item Draw the original circuit with $p_1$, $p_2$ as labelled rectangles (statement figure only), then redraw the circuit with your resistor / voltage source and save it as \file{Lab1/Ex09.asc}.
\item Use the current-controlled voltage source \cmd{h}: its controlling current is the current of a voltage source (name it in the value field, together with the transresistance 3).
\item Add the \cmd{.op} analysis (\textit{Simulate $\rightarrow$ Edit Simulation Cmd $\rightarrow$ DC op pnt}), run the simulation and read the node voltages and branch currents in the operating-point window.
\item Be careful with signs: LTspice reports the current of a voltage source as positive when it flows \emph{into} its positive terminal, and the current of a resistor as positive when it flows from its first pin to its second pin.
}

\comparetable{
\cmprow{$v$ (V)}
\cmprow{$i_x$ (A)}
\cmprow{$p_1$: type and power (W)}
\cmprow{$p_2$: type and power (W)}
}

\commentlines{2}

\subsection{Exercise 10: Resistive network: find a voltage}
Given the circuit in Figure~\ref{fig:lab1_ex10}, find the voltage $V$ across resistor R2. You can do this in any way, but remember to explain it in detail. Then simulate the circuit to check the result.

\labfig{lab1_ex10.png}{9cm}{Exercise 10: find the voltage $V$}{fig:lab1_ex10}

\theory{11}

\simulation{
\item Draw the circuit of Figure~\ref{fig:lab1_ex10} in LTspice (components \cmd{res}/\cmd{voltage}/\cmd{current}, ground \cmd{0}) and save it as \file{Lab1/Ex10.asc}.
\item Add the \cmd{.op} analysis (\textit{Simulate $\rightarrow$ Edit Simulation Cmd $\rightarrow$ DC op pnt}), run the simulation and read the node voltages and branch currents in the operating-point window.
\item Label node E as the ground (\texttt{0}) and the other nodes A to D.
\item Be careful with signs: LTspice reports the current of a voltage source as positive when it flows \emph{into} its positive terminal, and the current of a resistor as positive when it flows from its first pin to its second pin.
}

\comparetable{
\cmprow{$V$ (V)}
}

\commentlines{2}

\newpage
\section{Advanced exercises}
\subsection{Advanced Exercise 1: maximum power transfer}
In the circuit of Figure~\ref{fig:lab1_adv01}, a 12\,V source with a 100\,$\Omega$ series resistor R1 feeds node $x$; a 220\,$\Omega$ resistor R2 connects node $x$ to ground, and a load resistor $R_L$ is also connected between node $x$ and ground.
\labfig{lab1_adv01.png}{9cm}{Advanced Exercise 1: circuit with a variable load $R_L$}{fig:lab1_adv01}
\subsubsection{Theory calculation}\vspace{-3pt}
\notesblock
\begin{enumerate}[label=\alph*.]
  \item Find the Thevenin equivalent (voltage $V_{th}$ and resistance $R_{th}$) seen by $R_L$.\\ \dotfill
  \item Express the power $P_L$ dissipated in $R_L$ as a function of $R_L$, and determine the value of $R_L$ that maximises $P_L$ and the maximum power.\\ \dotfill
\end{enumerate}
\dotlines{4}
\simulation{
\item Define the load as a parameter: \cmd{.param RL=100}, and enter \cmd{\{RL\}} as the value of the load resistor.
\item Add \cmd{.step param RL 10 200 10} to simulate $R_L$ from 10 to 200\,$\Omega$ in steps of 10\,$\Omega$.
\item To evaluate one measurement per step, use a one-point DC sweep of the source: \cmd{.dc V1 0 12 12} (the sweep goes from 0\,V to 12\,V in one 12\,V step).
\item Measure the load power in every step with \cmd{.meas DC PL FIND V(x)*I(RL) AT=12}; the results appear in the SPICE Error Log. If your version rejects this expression, measure \cmd{V(x)} and \cmd{I(RL)} separately and compute the product yourself.
\item Fill in the table below (use the steps closest to the listed values) and find the load for which $P_L$ is maximum.
}
\begin{center}\renewcommand{\arraystretch}{1.7}
\begin{tabular}{|c|c|c|c|}\hline
\textbf{$R_L$ ($\Omega$)} & \textbf{$V_L$ (V)} & \textbf{$I_L$ (mA)} & \textbf{$P_L$ (mW)}\\ \hline
20 & & & \\ \hline 40 & & & \\ \hline 60 & & & \\ \hline 80 & & & \\ \hline 100 & & & \\ \hline 150 & & & \\ \hline 200 & & & \\ \hline
\end{tabular}\end{center}
\commentlines{2}

\newpage
\subsection{Advanced Exercise 2: superposition with \texttt{.step}}
In the circuit of Figure~\ref{fig:lab1_adv02}, two sources feed node $x$: $V_1=12$\,V through R1\,=\,1\,k$\Omega$ and $V_2=6$\,V through R2\,=\,2.2\,k$\Omega$; resistor R3\,=\,3.3\,k$\Omega$ connects node $x$ to ground.
\labfig{lab1_adv02.png}{9cm}{Advanced Exercise 2: two-source circuit}{fig:lab1_adv02}
\subsubsection{Theory calculation}\vspace{-3pt}
\notesblock
Find $V_x$ (a) with only $V_1$ active, (b) with only $V_2$ active, and (c) with both sources active. Show that (c) equals the sum of (a) and (b).
\dotlines{6}
\simulation{
\item Make the value of $V_1$ depend on a parameter: set it to \cmd{\{12*k1\}} and place \cmd{.param k1=1}.
\item Add \cmd{.step param k1 list 0 1} (source $V_1$ off, then on) and a one-point sweep of $V_2$: \cmd{.dc V2 0 6 6} (source $V_2$ off, then on).
\item Measure $V_x$ with \cmd{.meas DC Vx\_V2off FIND V(x) AT=0} and \cmd{.meas DC Vx\_V2on FIND V(x) AT=6}; read the four values (one per combination) in the SPICE Error Log.
}
\begin{center}\renewcommand{\arraystretch}{1.8}
\begin{tabular}{|p{5.2cm}|c|c|c|}\hline
\textbf{Case} & \textbf{Theory (V)} & \textbf{LTspice (V)} & \textbf{Error (\%)}\\ \hline
Only $V_1$ ($V_2$ off) & & & \\ \hline
Only $V_2$ ($V_1$ off) & & & \\ \hline
Both sources on & & & \\ \hline
Sum of the two single-source cases & & & \\ \hline
\end{tabular}\end{center}
\commentlines{2}

\section*{Table of standard resistor values}
\addcontentsline{toc}{section}{Table of standard resistor values (for Exercise 5)}
Standard values for resistors with 10\,\% tolerance (E12 series, in ohms; multiply by powers of ten):
\begin{center}\renewcommand{\arraystretch}{1.4}
\begin{tabular}{|*{12}{c|}}\hline
10 & 12 & 15 & 18 & 22 & 27 & 33 & 39 & 47 & 56 & 68 & 82\\ \hline
\end{tabular}\end{center}

\chap{Diode Circuits Analysis}

\section{Introduction}
In this laboratory you analyse circuits with diodes: the diode models, rectifiers, clampers, Zener regulators and a switching regulator. You compare the hand calculation (with the constant-voltage diode model) with the simulation (with a real diode model) and you discuss the differences.

\textbf{After this laboratory you will be able to:}
\begin{itemize}
  \item declare and use diode models, and plot the V-I characteristic with a DC sweep;
  \item analyse diode circuits with the constant-voltage model and verify them with \cmd{.op};
  \item run transient simulations of rectifiers and read the peak, ripple and conduction time;
  \item build and simulate a Zener regulator and a buck converter;
  \item use \cmd{.step} and \cmd{.meas} to study parameters and temperature.
\end{itemize}

\textbf{Exercises.} Save every schematic as \file{Lab2/ExNN.asc} (advanced exercises: \file{Lab2/AdvNN.asc}). Calculation spaces are blank; fill in the comparison tables with your own theoretical and simulated values.

\section{LTspice guide}

\subsection{Diode models}
Place a diode with the \texttt{D} key. A diode needs a model: either choose one from the library with \textit{Pick New Diode} in the component dialog, or declare it yourself with a SPICE directive (press \texttt{S}), for example \cmd{.model D1N4148 D(Is=2.52n Rs=0.568 N=1.752)}, and type the model name in the value field of the diode. A diode without parameters uses the default model of LTspice, whose parameters you can look up in the LTspice help (they are different from those of a real diode).

\labfig{lab2_guide01.png}{10cm}{Placing a diode and declaring its model}{fig:lab2_guide01}

\subsection{DC sweep and plotting}
Choose \textit{Simulate $\rightarrow$ Edit Simulation Cmd}, select the \textit{DC sweep} tab, enter the name of the source, the start value, the stop value and the increment. After the run, the waveform viewer opens; click a node to plot its voltage and click a component to plot its current. Click the name of a trace above the plot to attach a cursor.

\labfig{lab2_guide02.png}{12cm}{DC sweep settings and the waveform viewer}{fig:lab2_guide02}

\subsection{Transient analysis and sources}
Select the \textit{Transient} tab and enter the stop time and the maximum time step. A sine source is set from the voltage source dialog: click \textit{Advanced} and choose SINE (amplitude, frequency); a pulse train is set in the same way with PULSE (initial value, pulsed value, delay, rise and fall times, on time, period).

\labfig{lab2_guide03.png}{10cm}{Sine source and transient settings}{fig:lab2_guide03}

\subsection{Measurements, parameters and temperature}
The directive \cmd{.step} repeats a simulation for several values of a parameter or of the temperature (\cmd{.step param R list 1k 2k}, \cmd{.step temp list 25 75}); \cmd{.meas} measures a value automatically (for example \cmd{.meas DC Id FIND I(D1) AT=1}). Its results are written in the SPICE Error Log (\textit{View $\rightarrow$ SPICE Error Log}).

\subsection{Transformers}
LTspice has no ideal transformer: use two inductors and the directive \cmd{K1 L1 L2 1}. The turns ratio is the square root of the inductance ratio. Give the secondary a DC path to ground with a large resistor.

\labfig{lab2_guide05.png}{9cm}{A transformer made from two coupled inductors}{fig:lab2_guide05}

\section{Exercises}

\subsection{Exercise 1: Complete diode model}
According to the simulation circuit of Figure~\ref{fig:lab2_ex01} (a 5\,V source, a resistor R1 and one diode), determine the voltage across the resistor and the diode ($V_R$, $V_D$) and the current $I$ in the circuit with two different values of R1: 220\,$\Omega$ and 1.5\,k$\Omega$. Assume that the complete diode model is used for the calculation, with a forward voltage of 0.7\,V and an internal resistance of 50\,m$\Omega$. Finally, run the simulation to double-check the theory and explain briefly the difference between theory and simulation.

\labfig{lab2_ex01.png}{8cm}{Exercise 1: diode with a series resistor}{fig:lab2_ex01}

\theory{8}

\simulation{
\item Draw the circuit of Figure~\ref{fig:lab2_ex01} in LTspice and save it as \file{Lab2/Ex01.asc}.
\item Declare the diode model with a SPICE directive (\texttt{S} key), for example \cmd{.model D1N4148 D(Is=2.52n Rs=0.568 N=1.752)}, and type the model name \texttt{D1N4148} in the value field of every diode.
\item Define R1 as a parameter (\cmd{.param R=220}, value \cmd{\{R\}}) and run both cases with \cmd{.step param R list 220 1.5k}.
\item Add the \cmd{.op} analysis (\textit{Simulate $\rightarrow$ Edit Simulation Cmd $\rightarrow$ DC op pnt}), run the simulation and read the node voltages and branch currents.
\item Plot the V-I characteristic of the diode with a DC sweep: \cmd{.dc V1 0 5 10m}, and plot \cmd{I(D1)} against \cmd{V(a)}.
\item Be careful with signs: LTspice reports the current of a voltage source as positive when it flows \emph{into} its positive terminal, and the current of a two-terminal device as positive when it flows from its first pin to its second pin.
}

\comparetable{
\cmprow{$V_R$, R1\,=\,220\,$\Omega$ (V)}
\cmprow{$V_D$, R1\,=\,220\,$\Omega$ (V)}
\cmprow{$I$, R1\,=\,220\,$\Omega$ (mA)}
\cmprow{$V_R$, R1\,=\,1.5\,k$\Omega$ (V)}
\cmprow{$V_D$, R1\,=\,1.5\,k$\Omega$ (V)}
\cmprow{$I$, R1\,=\,1.5\,k$\Omega$ (mA)}
}

\commentlines{2}

\subsection{Exercise 2: Diodes in series}
Similar to the previous exercise, determine the voltages $V_{D1}$, $V_{D2}$, $V_{D3}$ and the current $I$ for the circuit of Figure~\ref{fig:lab2_ex02} (a 10\,V source, three diodes in series and a 1\,k$\Omega$ resistor). Use the practical diode model with a forward voltage of 0.7\,V in the calculation, then simulate the circuit again.

\labfig{lab2_ex02.png}{10cm}{Exercise 2: three diodes in series}{fig:lab2_ex02}

\theory{8}

\simulation{
\item Draw the circuit of Figure~\ref{fig:lab2_ex02} in LTspice and save it as \file{Lab2/Ex02.asc}.
\item Declare the diode model with a SPICE directive (\texttt{S} key), for example \cmd{.model D1N4148 D(Is=2.52n Rs=0.568 N=1.752)}, and type the model name \texttt{D1N4148} in the value field of every diode.
\item Add the \cmd{.op} analysis (\textit{Simulate $\rightarrow$ Edit Simulation Cmd $\rightarrow$ DC op pnt}), run the simulation and read the node voltages and branch currents.
\item Be careful with signs: LTspice reports the current of a voltage source as positive when it flows \emph{into} its positive terminal, and the current of a two-terminal device as positive when it flows from its first pin to its second pin.
}

\comparetable{
\cmprow{$V_{D1}$ (V)}
\cmprow{$V_{D2}$ (V)}
\cmprow{$V_{D3}$ (V)}
\cmprow{$V_{R1}$ (V)}
\cmprow{$I$ (mA)}
}

\commentlines{2}

\subsection{Exercise 3: Circuit analysis with a diode}
Consider the circuit of Figure~\ref{fig:lab2_ex03}: four resistors (R1 = 2\,k$\Omega$, R2 = 8\,k$\Omega$, R3 = 4\,k$\Omega$, R4 = 6\,k$\Omega$) and a diode D1 connected between the two middle nodes, supplied by a source of $V_s$. Assume that the diode is forward biased with a drop-down voltage of 0.7\,V. Write the nodal equations for the voltages $V_1$ (anode) and $V_2$ (cathode) and determine the current in every resistor for $V_s=10$\,V; solve it again for $V_s=8$\,V and 12\,V. Then run the simulation to check your results.

\labfig{lab2_ex03.png}{9cm}{Exercise 3: resistor network with a diode}{fig:lab2_ex03}

\theory{10}

\simulation{
\item Draw the circuit of Figure~\ref{fig:lab2_ex03} in LTspice and save it as \file{Lab2/Ex03.asc}.
\item Declare the diode model with a SPICE directive (\texttt{S} key), for example \cmd{.model D1N4148 D(Is=2.52n Rs=0.568 N=1.752)}, and type the model name \texttt{D1N4148} in the value field of every diode.
\item Define \cmd{.param Vs=10}, set the source to \cmd{\{Vs\}} and simulate the three cases with \cmd{.step param Vs list 8 10 12}.
\item Add the \cmd{.op} analysis (\textit{Simulate $\rightarrow$ Edit Simulation Cmd $\rightarrow$ DC op pnt}), run the simulation and read the node voltages and branch currents.
\item Be careful with signs: LTspice reports the current of a voltage source as positive when it flows \emph{into} its positive terminal, and the current of a two-terminal device as positive when it flows from its first pin to its second pin.
}

\comparetable{
\cmprow{$V_1$, $V_s=8$\,V (V)}
\cmprow{$V_2$, $V_s=8$\,V (V)}
\cmprow{$I_D$, $V_s=8$\,V (mA)}
\cmprow{$V_1$, $V_s=12$\,V (V)}
\cmprow{$V_2$, $V_s=12$\,V (V)}
\cmprow{$I_D$, $V_s=12$\,V (mA)}
}

\commentlines{2}

\subsection{Exercise 4: Clamper diode circuit}
The circuit of Figure~\ref{fig:lab2_ex04} (20\,V source V1, diodes D1 and D2, R1 = 5.6\,k$\Omega$, R2 = 3.3\,k$\Omega$) is used as a clamper, i.e.\ it clamps the voltage across R2 to a specific DC level. Using the practical diode model ($V_F=0.7$\,V), determine for $V_1=20$\,V which diodes conduct, the voltage across R2 and the currents $I_{R1}$, $I_{R2}$, $I_{D2}$. Then repeat the calculation for $V_1=12$\,V and explain the conclusion. Verify with a simulation.

\labfig{lab2_ex04.png}{9cm}{Exercise 4: clamper circuit}{fig:lab2_ex04}

\theory{10}

\simulation{
\item Draw the circuit of Figure~\ref{fig:lab2_ex04} in LTspice and save it as \file{Lab2/Ex04.asc}.
\item Declare the diode model with a SPICE directive (\texttt{S} key), for example \cmd{.model D1N4148 D(Is=2.52n Rs=0.568 N=1.752)}, and type the model name \texttt{D1N4148} in the value field of every diode.
\item Run both cases with \cmd{.param Vs=20} (value \cmd{\{Vs\}}) and \cmd{.step param Vs list 12 20}.
\item Add the \cmd{.op} analysis (\textit{Simulate $\rightarrow$ Edit Simulation Cmd $\rightarrow$ DC op pnt}), run the simulation and read the node voltages and branch currents.
\item Be careful with signs: LTspice reports the current of a voltage source as positive when it flows \emph{into} its positive terminal, and the current of a two-terminal device as positive when it flows from its first pin to its second pin.
}

\comparetable{
\cmprow{$I_{R1}$, 12\,V (mA)}
\cmprow{$I_{R2}$, 12\,V (mA)}
\cmprow{$I_{D2}$, 12\,V (mA)}
\cmprow{$V_{R2}$, 12\,V (V)}
\cmprow{$I_{R1}$, 20\,V (mA)}
\cmprow{$I_{R2}$, 20\,V (mA)}
\cmprow{$I_{D2}$, 20\,V (mA)}
\cmprow{$V_{R2}$, 20\,V (V)}
}

\commentlines{2}

\subsection{Exercise 5: Power switching circuit}
When power is abruptly lost, some systems switch automatically to a secondary source. In the circuit of Figure~\ref{fig:lab2_ex05} a 5\,V backup supply and a 9\,V main supply feed a 1\,k$\Omega$ load R$_L$, each through one diode. Assuming the practical diode model ($V_F=0.7$\,V), present the equations that give $I_{D3}$, $I_{D4}$, $I_{RL}$ and $V_{RL}$ in three cases: only the 5\,V source, only the 9\,V source, and both sources. Then simulate all cases.

\labfig{lab2_ex05.png}{9cm}{Exercise 5: power switching circuit}{fig:lab2_ex05}

\theory{10}

\simulation{
\item Draw the circuit of Figure~\ref{fig:lab2_ex05} in LTspice and save it as \file{Lab2/Ex05.asc}.
\item Declare the diode model with a SPICE directive (\texttt{S} key), for example \cmd{.model D1N4148 D(Is=2.52n Rs=0.568 N=1.752)}, and type the model name \texttt{D1N4148} in the value field of every diode.
\item Make the three cases with parameters: set the sources to \cmd{\{5*k3\}} and \cmd{\{9*k2\}}, and add \cmd{.step param k3 list 0 1} and \cmd{.step param k2 list 0 1}.
\item Add the \cmd{.op} analysis (\textit{Simulate $\rightarrow$ Edit Simulation Cmd $\rightarrow$ DC op pnt}), run the simulation and read the node voltages and branch currents.
}

\comparetable{
\cmprow{$I_{D3}$, 5\,V only (mA)}
\cmprow{$V_{RL}$, 5\,V only (V)}
\cmprow{$I_{D4}$, 9\,V only (mA)}
\cmprow{$V_{RL}$, 9\,V only (V)}
\cmprow{$I_{D3}$, both sources (mA)}
\cmprow{$I_{D4}$, both sources (mA)}
\cmprow{$V_{RL}$, both sources (V)}
}

\commentlines{2}

\subsection{Exercise 6: Half-wave rectifier}
In this exercise a sinusoidal source (amplitude 22\,V, frequency 50\,Hz) feeds a diode D1 and a load R1 = 1\,k$\Omega$ (see Figure~\ref{fig:lab2_ex06}). Using $V_F=0.78$\,V, calculate the maximum and the minimum value of the load voltage, the period of the source and the duration (conduction angle and time) during which the diode conducts in every cycle. Then run a transient simulation of 100\,ms and read the maximum and minimum values of the output.

\labfig{lab2_ex06.png}{8cm}{Exercise 6: half-wave rectifier}{fig:lab2_ex06}

\theory{12}

\simulation{
\item Draw the circuit of Figure~\ref{fig:lab2_ex06} in LTspice and save it as \file{Lab2/Ex06.asc}.
\item Declare the diode model with a SPICE directive (\texttt{S} key), for example \cmd{.model D1N4148 D(Is=2.52n Rs=0.568 N=1.752)}, and type the model name \texttt{D1N4148} in the value field of every diode.
\item Set the source to a sine wave: right-click V1 and use \textit{Advanced} to select SINE (amplitude 22\,V, frequency 50\,Hz).
\item Place a voltage marker on the output and use the transient analysis \cmd{.tran 0 100m 0 0.1m} (stop time 100\,ms, maximum time step 0.1\,ms).
\item Read the maximum and minimum with the cursors (click the trace name) or with \cmd{.meas TRAN Vmax MAX V(o)} and \cmd{.meas TRAN Vmin MIN V(o)}.
}

\comparetable{
\cmprow{$V_{R1,max}$ (V)}
\cmprow{$V_{R1,min}$ (V)}
\cmprow{Period $T$ (ms)}
\cmprow{Conduction time per cycle (ms)}
}

\commentlines{2}

\subsection{Exercise 7: Full-wave bridge rectifier}
The circuit of Figure~\ref{fig:lab2_ex07} is a full-wave bridge rectifier fed by a transformer with ratio $N_1/N_2=10$ from a 220\,V (peak), 50\,Hz source. Write the voltage equations $V_{AB}$ (transformer secondary) and $V_{CD}$ (load), assuming $V_F=0.7$\,V. Then perform a time-domain analysis to check them. Determine the time of the analysis for 10 periods of $V_{AB}$ and a time step that is ten times finer than the period of the signal.

\labfig{lab2_ex07.png}{11cm}{Exercise 7: full-wave bridge rectifier with a transformer}{fig:lab2_ex07}

\theory{10}

\simulation{
\item Draw the circuit of Figure~\ref{fig:lab2_ex07} in LTspice and save it as \file{Lab2/Ex07.asc}.
\item Declare the diode model with a SPICE directive (\texttt{S} key), for example \cmd{.model D1N4148 D(Is=2.52n Rs=0.568 N=1.752)}, and type the model name \texttt{D1N4148} in the value field of every diode.
\item Build the transformer with two inductors, for example L1\,=\,10\,H and L2\,=\,0.1\,H, and add the SPICE directive \cmd{K1 L1 L2 1}; the inductance ratio gives the turns ratio ($L_1/L_2=N_1^2/N_2^2$).
\item The secondary winding must have a DC path to ground: connect a large resistor (for example 100\,M$\Omega$) between one terminal of L2 and ground.
\item Plot $V_{AB}$ and $V_{CD}$ with a transient analysis (\cmd{.tran}) using your own stop time and maximum step.
}

\comparetable{
\cmprow{Amplitude of $V_{AB}$ (V)}
\cmprow{Peak of $V_{CD}$ (V)}
\cmprow{Period of $V_{AB}$ (ms)}
\cmprow{Stop time for 10 periods (ms)}
\cmprow{Maximum time step (ms)}
}

\commentlines{2}

\subsection{Exercise 8: Zener diode as a regulator}
A Zener diode is used to build the voltage regulator of Figure~\ref{fig:lab2_ex08} (R$_S$ = 200\,$\Omega$, R$_L$ = 1.5\,k$\Omega$, $V_Z=5$\,V). When the Zener diode operates in reverse breakdown, give the equations for the load current $I_L$, the series current $I_S$, the Zener current $I_Z$, the power $P_{R_S}$ dissipated in R$_S$ and the power $P_Z$ in the Zener diode. Then calculate them for $V_{CC}=8$\,V and 12\,V, and run the simulations (bias point) to confirm.

\labfig{lab2_ex08.png}{9cm}{Exercise 8: voltage regulator using a Zener diode}{fig:lab2_ex08}

\theory{10}

\simulation{
\item Draw the circuit of Figure~\ref{fig:lab2_ex08} in LTspice and save it as \file{Lab2/Ex08.asc}.
\item Declare the Zener model with a directive such as \cmd{.model DZ5V D(Is=1e-14 Bv=5.0 Ibv=10m Rs=2)} and use it in the diode that is connected in the reverse direction (cathode to the output node).
\item Simulate both input voltages with \cmd{.param Vcc=8} (value \cmd{\{Vcc\}}) and \cmd{.step param Vcc list 8 12}.
\item Add the \cmd{.op} analysis (\textit{Simulate $\rightarrow$ Edit Simulation Cmd $\rightarrow$ DC op pnt}), run the simulation and read the node voltages and branch currents.
\item The power in a component is the product of its voltage and current.
\item Be careful with signs: LTspice reports the current of a voltage source as positive when it flows \emph{into} its positive terminal, and the current of a two-terminal device as positive when it flows from its first pin to its second pin.
}

\comparetable{
\cmprow{$I_S$, 8\,V (mA)}
\cmprow{$I_L$, 8\,V (mA)}
\cmprow{$I_Z$, 8\,V (mA)}
\cmprow{$V_L$, 8\,V (V)}
\cmprow{$P_{R_S}$, 8\,V (mW)}
\cmprow{$P_Z$, 8\,V (mW)}
\cmprow{$I_S$, 12\,V (mA)}
\cmprow{$I_Z$, 12\,V (mA)}
\cmprow{$V_L$, 12\,V (V)}
\cmprow{$P_Z$, 12\,V (mW)}
}

\commentlines{2}

\subsection{Exercise 9: AC/DC power circuit application}
In this exercise you build, step by step, an AC-to-DC voltage source (Figure~\ref{fig:lab2_ex09}). The source is a 220\,V (peak), 50\,Hz sine wave, the transformer ratio is 10:1, the bridge uses four diodes and the load is R1 = 1\,k$\Omega$. Perform a time-domain simulation of 100\,ms for each step and write out comments and explanations.\begin{enumerate}[label=\textbf{Step \arabic*:},leftmargin=*,itemsep=2pt]\item The rectified voltage without any filtering or regulation.\item Add a 10\,$\micro$F filter capacitor in parallel with R1.\item Replace the capacitor by 680\,$\micro$F, re-run the simulation and explain the change.\item Add a Zener diode (reverse biased) in parallel with the load, with Zener voltage 22\,V; simulate and comment.\item Change the Zener voltage to 20\,V, re-run the simulation and explain any change.\end{enumerate}

\labfig{lab2_ex09.png}{12cm}{Exercise 9: bridge rectifier with filter capacitor and Zener diode (Step 5)}{fig:lab2_ex09}

\subsubsection{Theory calculation}\vspace{-3pt}\notesblock Estimate the peak output voltage of the rectifier and the ripple voltage expected for the 10\,$\micro$F and the 680\,$\micro$F capacitor (use $\Delta V\approx I/(2fC)$).\dotsx{2pt}{26.5pt}{6}

\simulation{
\item Draw the circuit of Figure~\ref{fig:lab2_ex09} in LTspice and save it as \file{Lab2/Ex09.asc}.
\item Declare the diode model with a SPICE directive (\texttt{S} key), for example \cmd{.model D1N4148 D(Is=2.52n Rs=0.568 N=1.752)}, and type the model name \texttt{D1N4148} in the value field of every diode.
\item Reuse the transformer and the bridge of Exercise~7. Modify the circuit step by step and save each version (for example \file{Lab2/Ex09\_step1.asc}, \dots).
\item Put two voltage markers: the transformer secondary and the output. Use \cmd{.tran} with a stop time of 100\,ms and a maximum step of 0.1\,ms.
\item For every step, measure the peak value and the ripple (difference between the maximum and the minimum of the output after the first 50\,ms) with the cursors or with \cmd{.meas TRAN Vpp PP V(out) FROM=50m TO=100m}.
}

\comparetable{
\cmprow{Step 1: output peak (V)}
\cmprow{Step 2: ripple (V)}
\cmprow{Step 3: ripple (V)}
\cmprow{Step 4: output level (V)}
\cmprow{Step 5: output level (V)}
}

\commentlines{6}

\subsection{Exercise 10: Buck converter (switching power supply)}
A switching regulator reduces a DC voltage efficiently. The circuit of Figure~\ref{fig:lab2_ex10} is a buck converter built from discrete components: a 12\,V source V1, a switch S1 driven by a 100\,kHz pulse train with a 50\,\% duty cycle, a diode D1, an inductor L1\,=\,33\,$\micro$H, a capacitor C1\,=\,220\,$\micro$F and a 10\,$\Omega$ load. Calculate the expected average output voltage and the load current, the peak-to-peak ripple of the inductor current $\Delta I_L=V_{out}(1-D)/(L f_{sw})$ and the output ripple voltage. Then run a transient simulation and compare. (The original exercise used the LM2596 regulator, whose vendor model cannot be assumed to be available; this circuit has the same learning objectives: switching regulation, inductor filtering and ripple.)

\labfig{lab2_ex10.png}{11cm}{Exercise 10: buck converter}{fig:lab2_ex10}

\theory{10}

\simulation{
\item Draw the circuit of Figure~\ref{fig:lab2_ex10} in LTspice and save it as \file{Lab2/Ex10.asc}.
\item Declare the switch with the directive \cmd{.model SW SW(Ron=0.05 Roff=1Meg Vt=2.5 Vh=0.1)} and place the switch (\cmd{sw}). Drive its control pins with a pulse source: \cmd{PULSE(0 5 0 10n 10n 5u 10u)}.
\item Declare a Schottky diode, for example \cmd{.model DSCH D(Is=1u N=1.05 Rs=0.05)}.
\item Run \cmd{.tran 0 3m 0 0.2u}. Look at the output voltage in the last 0.5\,ms (zoom in) to measure the average value and the ripple, and at the current in L1.
\item Optionally add \cmd{.meas TRAN Vavg AVG V(out) FROM=2.5m TO=3m} to obtain the average.
}

\comparetable{
\cmprow{Average output voltage (V)}
\cmprow{Load current (A)}
\cmprow{Inductor current ripple (A)}
\cmprow{Output voltage ripple (mV)}
}

\commentlines{2}

\newpage
\section{Advanced exercises}

\subsection{Advanced Exercise 1: Forward voltage versus temperature}
In the circuit of Figure~\ref{fig:lab2_adv01}, a 5\,V source feeds a 1\,k$\Omega$ resistor in series with a diode. Investigate how the forward voltage of the diode changes with the temperature.

\labfig{lab2_adv01.png}{8cm}{Advanced Exercise 1: diode and series resistor}{fig:lab2_adv01}

\subsubsection{Theory calculation}\vspace{-3pt}\notesblock Using the diode equation, explain why $V_D$ decreases when the temperature increases and estimate the slope in mV/$^\circ$C.\dotlines{4}

\simulation{
\item Add \cmd{.step temp list -25 25 75 125} to simulate four temperatures, and a one-point DC sweep \cmd{.dc V1 0 5 5}.
\item Measure the diode voltage with \cmd{.meas DC Vd FIND V(a) AT=5} and read the four results in the SPICE Error Log.
\item Fill in the table and compute the slope $\Delta V_D/\Delta T$.
}
\begin{center}\renewcommand{\arraystretch}{1.8}\begin{tabular}{|c|c|c|}\hline \textbf{Temperature ($^\circ$C)} & \textbf{$V_D$ (V)} & \textbf{$I_D$ (mA)}\\ \hline $-25$ & & \\ \hline 25 & & \\ \hline 75 & & \\ \hline 125 & & \\ \hline \end{tabular}\end{center}
\commentlines{2}
\newpage

\subsection{Advanced Exercise 2: Ripple of a rectifier versus the filter capacitor}
In the circuit of Figure~\ref{fig:lab2_adv02} a 22\,V (peak), 50\,Hz source feeds a bridge rectifier with a filter capacitor $C$ and a 1\,k$\Omega$ load. Study how the ripple voltage depends on $C$.

\labfig{lab2_adv02.png}{9cm}{Advanced Exercise 2: bridge rectifier with a filter capacitor}{fig:lab2_adv02}

\subsubsection{Theory calculation}\vspace{-3pt}\notesblock Estimate the ripple voltage $\Delta V\approx V_p/(2fR_LC)$ for $C=100$, 470, 1000 and 2200\,$\micro$F.\dotlines{4}

\simulation{
\item Set the capacitor to \cmd{\{C\}} with \cmd{.param C=470u} and add \cmd{.step param C list 100u 470u 1000u 2200u}.
\item Run \cmd{.tran 0 200m 0 0.1m} and measure the ripple in steady state with \cmd{.meas TRAN Vpp PP V(c,d) FROM=100m TO=200m}.
\item Read the four results in the SPICE Error Log and compare them with your estimates.
}
\begin{center}\renewcommand{\arraystretch}{1.8}\begin{tabular}{|c|c|c|c|}\hline \textbf{$C$ ($\micro$F)} & \textbf{Ripple, theory (V)} & \textbf{Ripple, LTspice (V)} & \textbf{Error (\%)}\\ \hline 100 & & & \\ \hline 470 & & & \\ \hline 1000 & & & \\ \hline 2200 & & & \\ \hline \end{tabular}\end{center}
\commentlines{2}
\newpage

\chap{Bipolar Junction Transistor}

\section{Introduction}
In this laboratory you study the bipolar junction transistor (BJT): its characteristics, its use as a switch and as a driver, the common bias configurations, the Darlington pair, the common-base amplifier, the current mirror and two typical applications (a logic gate and an optocoupler).

\textbf{After this laboratory you will be able to:}
\begin{itemize}
  \item declare NPN and PNP models and plot the output characteristics with a nested DC sweep;
  \item decide the operating region of a BJT (active or saturated) from the simulated voltages;
  \item calculate and verify bias circuits (fixed, voltage-divider, emitter resistor, PNP);
  \item analyse a Darlington pair, a common-base stage and a current mirror;
  \item use a BJT in switching applications (switch, LED driver, NOT gate) and model an optocoupler;
  \item run an AC analysis to obtain the gain and the cut-off frequencies of an amplifier.
\end{itemize}

\textbf{Exercises.} Save every schematic as \file{Lab3/ExNN.asc} (advanced exercises: \file{Lab3/AdvNN.asc}). Calculation spaces are blank; fill in the comparison tables with your own theoretical and simulated values.

\section{LTspice guide}

\subsection{Bipolar transistors in LTspice}
Place a transistor from the component list (\texttt{F2}): \cmd{npn} or \cmd{pnp}. Use \textit{Pick New Transistor} in the dialog to select a library device, or declare the models with directives as shown in the exercises and type the model name in the value field. The arrow on the emitter shows the type; use \texttt{Ctrl+R} and \texttt{Ctrl+E} to rotate and mirror the symbol so that the terminals match your drawing.

\labfig{lab3_guide01.png}{10cm}{Placing an NPN transistor and giving it a model}{fig:lab3_guide01}

\subsection{Nested DC sweep}
To plot a family of output characteristics, sweep two sources: in \textit{Edit Simulation Cmd $\rightarrow$ DC sweep} fill in the first source (the collector voltage) and enable the second source tab (the base current), with start, stop and increment for each. The corresponding directive is \cmd{.dc V1 0 10 0.05 I1 20u 100u 20u}.

\labfig{lab3_guide02.png}{12cm}{Nested DC sweep: collector voltage and base current}{fig:lab3_guide02}

\subsection{Parameters, steps and the AC analysis}
A parameter is defined with \cmd{.param name=value} and used in curly brackets (\cmd{\{name\}}); it can be stepped with \cmd{.step param name list v1 v2 v3}. For a frequency response choose the \textit{AC Analysis} tab (decade sweep, points per decade, start and stop frequency) and give the input source an AC amplitude of 1. In the waveform viewer, the magnitude is displayed in dB together with the phase.

\labfig{lab3_guide03.png}{10cm}{AC analysis settings}{fig:lab3_guide03}

\section{Exercises}

\subsection{Exercise 1: Output characteristics and saturation}
Consider the NPN transistor circuits of Figure~\ref{fig:lab3_ex01} (a) and (b). In (a), V$_{CC}$ = 10\,V, R$_C$ = 1\,k$\Omega$, R$_B$ = 10\,k$\Omega$ and the base supply is 5\,V. Calculate $I_B$, $I_C$ and $V_{CE}$ and decide in which mode the transistor works. Assume $\beta=100$, $V_{BE}=0.7$\,V and $|V_{CE(sat)}|\approx0.2$\,V unless stated otherwise. In (b), the transistor is driven by a base current source and the collector voltage is swept. Plot the output characteristics $I_C=f(V_{CE})$ for $I_B=20$, 40, 60, 80 and 100\,$\micro$A, identify the saturation region and the active region, and estimate $\beta$ from the curves.

\labfig{lab3_ex01.png}{8cm}{Exercise 1 (a): transistor in saturation}{fig:lab3_ex01}

\labfig{lab3_ex01b.png}{8cm}{Exercise 1 (b): circuit for the output characteristics}{fig:lab3_ex01b}

\theory{8}

\simulation{
\item Draw the circuit of Figure~\ref{fig:lab3_ex01} in LTspice and save it as \file{Lab3/Ex01.asc}.
\item Declare the transistor models with SPICE directives (\texttt{S} key): \cmd{.model QN NPN(Is=1e-14 Bf=100 Br=1)} for an NPN transistor and \cmd{.model QP PNP(Is=1e-14 Bf=100 Br=1)} for a PNP transistor. Place the transistor (\cmd{npn} or \cmd{pnp}) and type the model name in its value field.
\item Circuit (a): run \cmd{.op} and read $V_{CE}$ and $I_C$.
\item Circuit (b): use a nested DC sweep \cmd{.dc V1 0 10 0.05 I1 20u 100u 20u} (first source: the collector voltage; second source: the base current) and plot \cmd{Ic(Q1)}.
\item Be careful with signs: LTspice reports the current of a voltage source as positive when it flows \emph{into} its positive terminal, and the current of a two-terminal device as positive when it flows from its first pin to its second pin.
}

\comparetable{
\cmprow{(a) $I_B$ (mA)}
\cmprow{(a) $I_C$ (mA)}
\cmprow{(a) $V_{CE}$ (V)}
\cmprow{(b) $\beta$ at $V_{CE}=5$\,V ($I_B=40\,\micro$A)}
}

\commentlines{2}

\subsection{Exercise 2: BJT as a switch and as a device driver}
(a) For the BJT switch of Figure~\ref{fig:lab3_ex02}(a) (V$_{CC}$ = 5\,V, input 5\,V) determine R$_C$ and R$_B$ such that the collector current saturates at 50\,mA, with $V_{CE(sat)}=0.2$\,V, $V_{BE}=0.7$\,V and a forced current gain $\beta_{forced}=10$. (b) A red LED has to carry 10\,mA from a 5\,V supply, switched by a 3.3\,V logic signal (Figure~\ref{fig:lab3_ex02b}): determine the series resistor of the LED (LED voltage about 1.8\,V) and the base resistor. Use standard resistor values and simulate both circuits with them.

\labfig{lab3_ex02.png}{8cm}{Exercise 2 (a): BJT switch in saturation}{fig:lab3_ex02}

\labfig{lab3_ex02b.png}{8cm}{Exercise 2 (b): driving an LED with a BJT}{fig:lab3_ex02b}

\theory{10}

\simulation{
\item Draw the circuit of Figure~\ref{fig:lab3_ex02} in LTspice and save it as \file{Lab3/Ex02.asc}.
\item Declare the transistor models with SPICE directives (\texttt{S} key): \cmd{.model QN NPN(Is=1e-14 Bf=100 Br=1)} for an NPN transistor and \cmd{.model QP PNP(Is=1e-14 Bf=100 Br=1)} for a PNP transistor. Place the transistor (\cmd{npn} or \cmd{pnp}) and type the model name in its value field.
\item Declare an LED model with a directive such as \cmd{.model LEDRED D(Is=1e-18 N=1.9 Rs=5)}.
\item Add the \cmd{.op} analysis (\textit{Simulate $\rightarrow$ Edit Simulation Cmd $\rightarrow$ DC op pnt}), run the simulation and read the node voltages and branch currents.
\item Check that $V_{CE}$ is close to the saturation value and compare $I_C$ and the LED current with the specification.
\item Be careful with signs: LTspice reports the current of a voltage source as positive when it flows \emph{into} its positive terminal, and the current of a two-terminal device as positive when it flows from its first pin to its second pin.
}

\comparetable{
\cmprow{(a) $R_C$ calculated ($\Omega$)}
\cmprow{(a) $R_B$ calculated ($\Omega$)}
\cmprow{(a) $I_C$ with standard values (mA)}
\cmprow{(a) $V_{CE}$ (V)}
\cmprow{(b) LED current (mA)}
\cmprow{(b) $V_{CE}$ (V)}
}

\commentlines{2}

\subsection{Exercise 3: Simple bias: voltage-divider configuration}
In the circuit of Figure~\ref{fig:lab3_ex03} (V$_{CC}$ = 12\,V, R1 = 47\,k$\Omega$, R2 = 10\,k$\Omega$, R$_C$ = 2.2\,k$\Omega$, R$_E$ = 1\,k$\Omega$) determine the base voltage $V_B$, the emitter current $I_E$, the collector current $I_C$ and the voltages $V_C$ and $V_{CE}$, using the approximate method of the voltage divider (Assume $\beta=100$, $V_{BE}=0.7$\,V and $|V_{CE(sat)}|\approx0.2$\,V unless stated otherwise.). Then simulate and discuss the effect of the base current on the divider.

\labfig{lab3_ex03.png}{8cm}{Exercise 3: voltage-divider bias}{fig:lab3_ex03}

\theory{10}

\simulation{
\item Draw the circuit of Figure~\ref{fig:lab3_ex03} in LTspice and save it as \file{Lab3/Ex03.asc}.
\item Declare the transistor models with SPICE directives (\texttt{S} key): \cmd{.model QN NPN(Is=1e-14 Bf=100 Br=1)} for an NPN transistor and \cmd{.model QP PNP(Is=1e-14 Bf=100 Br=1)} for a PNP transistor. Place the transistor (\cmd{npn} or \cmd{pnp}) and type the model name in its value field.
\item Add the \cmd{.op} analysis (\textit{Simulate $\rightarrow$ Edit Simulation Cmd $\rightarrow$ DC op pnt}), run the simulation and read the node voltages and branch currents.
\item Be careful with signs: LTspice reports the current of a voltage source as positive when it flows \emph{into} its positive terminal, and the current of a two-terminal device as positive when it flows from its first pin to its second pin.
}

\comparetable{
\cmprow{$V_B$ (V)}
\cmprow{$V_E$ (V)}
\cmprow{$I_C$ (mA)}
\cmprow{$V_C$ (V)}
\cmprow{$V_{CE}$ (V)}
}

\commentlines{2}

\subsection{Exercise 4: PNP circuit and complementary pair}
(a) In the PNP circuit of Figure~\ref{fig:lab3_ex04} (supply +10\,V, R1 = 22\,k$\Omega$, R2 = 10\,k$\Omega$, R$_E$ = 2\,k$\Omega$ to the supply, R$_C$ = 680\,$\Omega$ to ground) determine $V_B$, $V_E$, $I_C$, $V_C$ and $V_{EC}$. Assume $\beta=100$, $V_{BE}=0.7$\,V and $|V_{CE(sat)}|\approx0.2$\,V unless stated otherwise. (b) The circuit of Figure~\ref{fig:lab3_ex04b} uses an NPN and a PNP transistor with supplies of $\pm12$\,V and a 100\,$\Omega$ load (complementary emitter follower). For a 5\,V peak, 1\,kHz sine input, find the expected output peak value and explain what happens when the input crosses zero.

\labfig{lab3_ex04.png}{8cm}{Exercise 4 (a): PNP bias circuit}{fig:lab3_ex04}

\labfig{lab3_ex04b.png}{9cm}{Exercise 4 (b): complementary emitter follower}{fig:lab3_ex04b}

\theory{10}

\simulation{
\item Draw the circuit of Figure~\ref{fig:lab3_ex04} in LTspice and save it as \file{Lab3/Ex04.asc}.
\item Declare the transistor models with SPICE directives (\texttt{S} key): \cmd{.model QN NPN(Is=1e-14 Bf=100 Br=1)} for an NPN transistor and \cmd{.model QP PNP(Is=1e-14 Bf=100 Br=1)} for a PNP transistor. Place the transistor (\cmd{npn} or \cmd{pnp}) and type the model name in its value field.
\item Circuit (a): run \cmd{.op}. Circuit (b): run \cmd{.tran 0 5m 0 5u} with a sine source (amplitude 5\,V, 1\,kHz) and plot the input and the output.
\item Zoom in around the zero crossings of the output and measure the dead zone.
\item Be careful with signs: LTspice reports the current of a voltage source as positive when it flows \emph{into} its positive terminal, and the current of a two-terminal device as positive when it flows from its first pin to its second pin.
}

\comparetable{
\cmprow{(a) $V_B$ (V)}
\cmprow{(a) $I_C$ (mA)}
\cmprow{(a) $V_C$ (V)}
\cmprow{(a) $V_{EC}$ (V)}
\cmprow{(b) output peak (V)}
\cmprow{(b) dead zone around zero (V)}
}

\commentlines{2}

\subsection{Exercise 5: NPN circuit with an emitter resistance}
For the circuit of Figure~\ref{fig:lab3_ex05} (V$_{CC}$ = 12\,V, R$_B$ = 270\,k$\Omega$, R$_C$ = 1.5\,k$\Omega$, R$_E$ = 1\,k$\Omega$) write the base-emitter loop equation, determine $I_B$, $I_C$, $V_E$, $V_C$ and $V_{CE}$, and verify that the transistor is in the active region. Assume $\beta=100$, $V_{BE}=0.7$\,V and $|V_{CE(sat)}|\approx0.2$\,V unless stated otherwise. Simulate and compare.

\labfig{lab3_ex05.png}{8cm}{Exercise 5: NPN with an emitter resistor}{fig:lab3_ex05}

\theory{10}

\simulation{
\item Draw the circuit of Figure~\ref{fig:lab3_ex05} in LTspice and save it as \file{Lab3/Ex05.asc}.
\item Declare the transistor models with SPICE directives (\texttt{S} key): \cmd{.model QN NPN(Is=1e-14 Bf=100 Br=1)} for an NPN transistor and \cmd{.model QP PNP(Is=1e-14 Bf=100 Br=1)} for a PNP transistor. Place the transistor (\cmd{npn} or \cmd{pnp}) and type the model name in its value field.
\item Add the \cmd{.op} analysis (\textit{Simulate $\rightarrow$ Edit Simulation Cmd $\rightarrow$ DC op pnt}), run the simulation and read the node voltages and branch currents.
\item Be careful with signs: LTspice reports the current of a voltage source as positive when it flows \emph{into} its positive terminal, and the current of a two-terminal device as positive when it flows from its first pin to its second pin.
}

\comparetable{
\cmprow{$I_B$ ($\micro$A)}
\cmprow{$I_C$ (mA)}
\cmprow{$V_E$ (V)}
\cmprow{$V_C$ (V)}
\cmprow{$V_{CE}$ (V)}
}

\commentlines{2}

\subsection{Exercise 6: Darlington circuit}
In the Darlington circuit of Figure~\ref{fig:lab3_ex06} a 3\,V source drives the base of Q1 through 470\,k$\Omega$; R2 = 4.7\,k$\Omega$ and R3 = 470\,$\Omega$ are connected to a 10\,V supply. Determine $I_{BE}$ (the current in the 470\,k$\Omega$ resistor), the current $I_{AC}$ through the 4.7\,k$\Omega$ resistor and the current $I_{AL}$ through the 470\,$\Omega$ resistor, and the overall current gain $I_{AL}/I_{BE}$. Assume both transistors have $\beta=100$ and check whether Q2 saturates.

\labfig{lab3_ex06.png}{8cm}{Exercise 6: Darlington circuit}{fig:lab3_ex06}

\theory{10}

\simulation{
\item Draw the circuit of Figure~\ref{fig:lab3_ex06} in LTspice and save it as \file{Lab3/Ex06.asc}.
\item Declare the transistor models with SPICE directives (\texttt{S} key): \cmd{.model QN NPN(Is=1e-14 Bf=100 Br=1)} for an NPN transistor and \cmd{.model QP PNP(Is=1e-14 Bf=100 Br=1)} for a PNP transistor. Place the transistor (\cmd{npn} or \cmd{pnp}) and type the model name in its value field.
\item Add the \cmd{.op} analysis (\textit{Simulate $\rightarrow$ Edit Simulation Cmd $\rightarrow$ DC op pnt}), run the simulation and read the node voltages and branch currents.
\item Compare the simulated overall gain with $\beta^2$ and explain the difference.
\item Be careful with signs: LTspice reports the current of a voltage source as positive when it flows \emph{into} its positive terminal, and the current of a two-terminal device as positive when it flows from its first pin to its second pin.
}

\comparetable{
\cmprow{$I_{BE}$ ($\micro$A)}
\cmprow{$I_{AC}$ (mA)}
\cmprow{$I_{AL}$ (mA)}
\cmprow{$I_{AL}/I_{BE}$}
\cmprow{Mode of Q2}
}

\commentlines{2}

\subsection{Exercise 7: Common-base circuit}
Figure~\ref{fig:lab3_ex07} shows a common-base circuit (R1 = 1.2\,k$\Omega$, V1 = 4\,V, R2 = 2.4\,k$\Omega$, V2 = 10\,V; the base is grounded and the emitter is returned to $-4$\,V). Calculate $I_E$, $I_B$, $I_C$, $V_C$ and $V_{CE}$. Assume $\beta=100$, $V_{BE}=0.7$\,V and $|V_{CE(sat)}|\approx0.2$\,V unless stated otherwise. Then simulate the circuit. Check that the transistor is wired in the right direction (the emitter arrow points away from the base).

\labfig{lab3_ex07.png}{9cm}{Exercise 7: common-base circuit}{fig:lab3_ex07}

\theory{10}

\simulation{
\item Draw the circuit of Figure~\ref{fig:lab3_ex07} in LTspice and save it as \file{Lab3/Ex07.asc}.
\item Declare the transistor models with SPICE directives (\texttt{S} key): \cmd{.model QN NPN(Is=1e-14 Bf=100 Br=1)} for an NPN transistor and \cmd{.model QP PNP(Is=1e-14 Bf=100 Br=1)} for a PNP transistor. Place the transistor (\cmd{npn} or \cmd{pnp}) and type the model name in its value field.
\item If $I_B$ is as large as $I_C$ in the simulation, the transistor is upside down: mirror it (\texttt{Ctrl+E}).
\item Add the \cmd{.op} analysis (\textit{Simulate $\rightarrow$ Edit Simulation Cmd $\rightarrow$ DC op pnt}), run the simulation and read the node voltages and branch currents.
\item Be careful with signs: LTspice reports the current of a voltage source as positive when it flows \emph{into} its positive terminal, and the current of a two-terminal device as positive when it flows from its first pin to its second pin.
}

\comparetable{
\cmprow{$I_E$ (mA)}
\cmprow{$I_B$ ($\micro$A)}
\cmprow{$I_C$ (mA)}
\cmprow{$V_C$ (V)}
\cmprow{$V_{CE}$ (V)}
}

\commentlines{2}

\subsection{Exercise 8: Current mirror}
The circuit of Figure~\ref{fig:lab3_ex08} is a current mirror (V1 = 12\,V, R2 = 1\,k$\Omega$ as the load, Q1 and Q2 identical, $\beta=100$). First calculate, for R1 = 2\,k$\Omega$, the control current $I_{CTRL}$, the base currents and the collector currents $I_{C1}$ and $I_L=I_{C2}$ and explain why the circuit is called a current mirror. Now replace R1 by 100\,$\Omega$, calculate all the values again, simulate the new circuit and explain the phenomenon you observe.

\labfig{lab3_ex08.png}{8cm}{Exercise 8: current mirror}{fig:lab3_ex08}

\theory{10}

\simulation{
\item Draw the circuit of Figure~\ref{fig:lab3_ex08} in LTspice and save it as \file{Lab3/Ex08.asc}.
\item Declare the transistor models with SPICE directives (\texttt{S} key): \cmd{.model QN NPN(Is=1e-14 Bf=100 Br=1)} for an NPN transistor and \cmd{.model QP PNP(Is=1e-14 Bf=100 Br=1)} for a PNP transistor. Place the transistor (\cmd{npn} or \cmd{pnp}) and type the model name in its value field.
\item Define R1 as a parameter and run both values with \cmd{.step param R1v list 2k 100} (R1 = \cmd{\{R1v\}}).
\item Add the \cmd{.op} analysis (\textit{Simulate $\rightarrow$ Edit Simulation Cmd $\rightarrow$ DC op pnt}), run the simulation and read the node voltages and branch currents.
\item Be careful with signs: LTspice reports the current of a voltage source as positive when it flows \emph{into} its positive terminal, and the current of a two-terminal device as positive when it flows from its first pin to its second pin.
}

\comparetable{
\cmprow{$I_{CTRL}$, R1\,=\,2\,k$\Omega$ (mA)}
\cmprow{$I_{C1}$, R1\,=\,2\,k$\Omega$ (mA)}
\cmprow{$I_L$, R1\,=\,2\,k$\Omega$ (mA)}
\cmprow{$I_{CTRL}$, R1\,=\,100\,$\Omega$ (mA)}
\cmprow{$I_{C1}$, R1\,=\,100\,$\Omega$ (mA)}
\cmprow{$I_L$, R1\,=\,100\,$\Omega$ (mA)}
}

\commentlines{2}

\subsection{Exercise 9: BJT logic gate: the NOT gate}
The circuit of Figure~\ref{fig:lab3_ex09} is a NOT gate (V2 = 5\,V, R$_C$ = 10\,k$\Omega$, R$_B$ = 10\,k$\Omega$). The input source V1 is a pulse source: V1 = 0\,V, V2 = 5\,V, TD = 0, TR = 5\,ns, TF = 3\,ns, PW = 50\,ms, PER = 100\,ms. Explain, with calculations, why the output is high for a low input and low for a high input (check the saturation of the transistor). Find the maximum current this gate can source and explain why adding an additional 220\,$\Omega$ resistor in parallel with R$_C$ allows the gate to drive a LED. Simulate the circuit with a transient analysis of 200\,ms.

\labfig{lab3_ex09.png}{8cm}{Exercise 9: BJT NOT gate}{fig:lab3_ex09}

\theory{10}

\simulation{
\item Draw the circuit of Figure~\ref{fig:lab3_ex09} in LTspice and save it as \file{Lab3/Ex09.asc}.
\item Declare the transistor models with SPICE directives (\texttt{S} key): \cmd{.model QN NPN(Is=1e-14 Bf=100 Br=1)} for an NPN transistor and \cmd{.model QP PNP(Is=1e-14 Bf=100 Br=1)} for a PNP transistor. Place the transistor (\cmd{npn} or \cmd{pnp}) and type the model name in its value field.
\item Set V1 as a pulse source (voltage source dialog $\rightarrow$ \textit{Advanced} $\rightarrow$ PULSE) with the parameters above.
\item Run \cmd{.tran 0 200m 0 0.1m} and plot the input and the collector voltage.
\item Optionally run a DC sweep of the input from 0 to 5\,V to obtain the transfer characteristic.
}

\comparetable{
\cmprow{Output for input 0\,V (V)}
\cmprow{Output for input 5\,V (V)}
\cmprow{$I_B$ for input 5\,V (mA)}
\cmprow{Maximum output current (mA)}
}

\commentlines{2}

\subsection{Exercise 10: Optocoupler}
In the circuit of Figure~\ref{fig:lab3_ex10} a 9\,V battery drives, through a 10\,k$\Omega$ resistor, the LED of an optocoupler whose phototransistor is pulled up to 5\,V by R2 = 4.7\,k$\Omega$ (output V$_{OUT}$). The LED forward voltage is 1.7\,V and the current transfer ratio (CTR) is 2, i.e.\ $I_C=CTR\cdot I_F$. Calculate the LED current $I_F$, the collector current the phototransistor would need, the current $I_{R2}$ and the output voltage V$_{OUT}$ when the switch is closed and when it is open; check whether the output transistor saturates. Explain what the optocoupler is used for. (The vendor model of the PC817 is not available in LTspice; the optocoupler is modelled by a LED, a current-controlled current source with gain CTR/$\beta$ driving the base of an NPN transistor.)

\labfig{lab3_ex10.png}{12cm}{Exercise 10: optocoupler circuit with a behavioural model}{fig:lab3_ex10}

\theory{10}

\simulation{
\item Draw the circuit of Figure~\ref{fig:lab3_ex10} in LTspice and save it as \file{Lab3/Ex10.asc}.
\item Declare the transistor models with SPICE directives (\texttt{S} key): \cmd{.model QN NPN(Is=1e-14 Bf=100 Br=1)} for an NPN transistor and \cmd{.model QP PNP(Is=1e-14 Bf=100 Br=1)} for a PNP transistor. Place the transistor (\cmd{npn} or \cmd{pnp}) and type the model name in its value field.
\item Use a zero-volt source Vs in series with the LED to sense its current. Place the current-controlled current source \cmd{f}: its value field contains the name of the sensing source and the gain 0.02 (= CTR/$\beta$ = 2/100); its output is connected from ground to the base of the output transistor.
\item Add the \cmd{.op} analysis (\textit{Simulate $\rightarrow$ Edit Simulation Cmd $\rightarrow$ DC op pnt}), run the simulation and read the node voltages and branch currents.
\item Open the switch (remove the source or set 0\,V) and run again.
\item Be careful with signs: LTspice reports the current of a voltage source as positive when it flows \emph{into} its positive terminal, and the current of a two-terminal device as positive when it flows from its first pin to its second pin.
}

\comparetable{
\cmprow{$I_F$ (mA)}
\cmprow{$I_{R2}$, switch closed (mA)}
\cmprow{$V_{OUT}$, switch closed (V)}
\cmprow{$V_{OUT}$, switch open (V)}
}

\commentlines{2}

\newpage
\section{Advanced exercises}

\subsection{Advanced Exercise 1: Q-point versus transistor gain}
Use the voltage-divider bias circuit of Exercise~3 (Figure~\ref{fig:lab3_adv01}). Study how the operating point changes when the current gain $\beta$ of the transistor changes from 50 to 200 (a typical production spread).

\labfig{lab3_adv01.png}{8cm}{Advanced Exercise 1: voltage-divider bias with a variable $\beta$}{fig:lab3_adv01}

\subsubsection{Theory calculation}\vspace{-3pt}\notesblock Calculate $I_C$ and $V_{CE}$ for $\beta=50$, 100 and 200 with the approximate method, and explain why this bias circuit is called stable.\dotlines{4}

\simulation{
\item Declare the model with a parameter: \cmd{.model QV NPN(Is=1e-14 Bf=\{BF\} Br=1)}, define \cmd{.param BF=100} and add \cmd{.step param BF list 50 100 200}.
\item Use a one-point DC sweep \cmd{.dc V1 12 12 1} and measure with \cmd{.meas DC Ic FIND I(RC) AT=12} and \cmd{.meas DC Vce FIND V(c,e) AT=12}.
\item Compare the variation of $I_C$ (in percent) with the variation of $\beta$.
}
\begin{center}\renewcommand{\arraystretch}{1.8}\begin{tabular}{|c|c|c|c|c|}\hline \textbf{$\beta$} & \textbf{$I_C$ theory (mA)} & \textbf{$I_C$ LTspice (mA)} & \textbf{$V_{CE}$ theory (V)} & \textbf{$V_{CE}$ LTspice (V)}\\ \hline 50 & & & & \\ \hline 100 & & & & \\ \hline 200 & & & & \\ \hline \end{tabular}\end{center}
\commentlines{2}
\newpage

\subsection{Advanced Exercise 2: Frequency response of a common-emitter amplifier}
Figure~\ref{fig:lab3_adv02} shows a common-emitter amplifier built from the bias circuit of Exercise~3, with a 10\,$\micro$F input capacitor, a 100\,$\micro$F emitter bypass capacitor, a 10\,$\micro$F output capacitor and a 10\,k$\Omega$ load. Determine the mid-band voltage gain and the lower and upper cut-off frequencies.

\labfig{lab3_adv02.png}{11cm}{Advanced Exercise 2: common-emitter amplifier}{fig:lab3_adv02}

\subsubsection{Theory calculation}\vspace{-3pt}\notesblock Estimate the small-signal mid-band gain $A_v\approx-g_m(R_C\parallel R_L)$, with $g_m=I_C/V_T$ ($V_T=25.9$\,mV).\dotlines{4}

\simulation{
\item Use a transistor model with capacitances, for example \cmd{.model QA NPN(Is=1e-14 Bf=100 Br=1 Vaf=100 Cje=10p Cjc=4p Tf=0.4n)}.
\item Set the AC amplitude of the input source to 1\,V (voltage source dialog) and run \cmd{.ac dec 20 10 100Meg}.
\item Plot \cmd{V(out)} and read the mid-band gain and the $-3$\,dB frequencies with the cursors.
}
\begin{center}\renewcommand{\arraystretch}{1.8}\begin{tabular}{|c|c|c|c|}\hline \textbf{Quantity} & \textbf{Theory} & \textbf{LTspice} & \textbf{Error (\%)}\\ \hline Mid-band gain (V/V) & & & \\ \hline Mid-band gain (dB) & & & \\ \hline Lower cut-off $f_L$ (Hz) & & & \\ \hline Upper cut-off $f_H$ (MHz) & & & \\ \hline \end{tabular}\end{center}
\commentlines{2}
\newpage

\chap{Operational Amplifier}

\section{Introduction}
In this laboratory you study circuits built with an operational amplifier: the basic amplifier configurations, the voltage follower and its current-boosted versions, summing and difference amplifiers, first-order active filters and the Schmitt trigger. You verify the ideal-op-amp equations with DC, AC and transient simulations.

\textbf{After this laboratory you will be able to:}
\begin{itemize}
  \item place an op-amp and its supply sources and use the virtual-short concept;
  \item calculate and simulate inverting, non-inverting, summing and difference amplifiers;
  \item use a BJT at the output of an op-amp to increase the output current;
  \item obtain the Bode plot of active low-pass and high-pass filters with an AC analysis;
  \item explain and measure the hysteresis of a Schmitt trigger in a transient simulation;
  \item observe the effect of the finite gain-bandwidth product of a real op-amp.
\end{itemize}

\textbf{Exercises.} Save every schematic as \file{Lab4/ExNN.asc} (advanced exercises: \file{Lab4/AdvNN.asc}). Calculation spaces are blank; fill in the comparison tables with your own theoretical and simulated values.

\section{LTspice guide}

\subsection{Op-amp symbols and supplies}
Place the op-amp symbol with \texttt{F2}: search for \cmd{UniversalOpAmp2} in the \textit{[Opamps]} folder. The symbol has five pins: non-inverting input (+), inverting input ($-$), output, and two supply pins. Connect a $+15$\,V source to the upper supply pin and a $-15$\,V source to the lower one; ground the other terminals of both sources. The output of the op-amp cannot exceed the supply rails: it saturates a little below them.

\labfig{lab4_guide01.png}{10cm}{An op-amp with its two supply sources}{fig:lab4_guide01}

\subsection{AC analysis and Bode plot}
Give the input source an AC amplitude of 1\,V (voltage source dialog, \textit{Small signal AC analysis}). Choose \textit{Simulate $\rightarrow$ Edit Simulation Cmd $\rightarrow$ AC Analysis}, select the decade sweep, the number of points per decade and the start and stop frequencies. Plot the output node: LTspice shows the magnitude in dB and the phase. The cut-off frequency is where the magnitude is 3\,dB below its pass-band value.

\labfig{lab4_guide02.png}{12cm}{AC analysis settings and Bode plot}{fig:lab4_guide02}

\subsection{Transient analysis of comparators}
A Schmitt trigger is studied in the time domain. Use a sine source and a \cmd{.tran} analysis; plot the input and the output in the same plot window (select the second node after the first). Use the cursors to read the input value at the instants when the output switches.

\section{Exercises}

\subsection{Exercise 1: Non-inverting amplifier}
In the circuit of Figure~\ref{fig:lab4_ex01} a 2\,V DC source drives the non-inverting input of an op-amp; the output is divided by R$_f$ = 1\,k$\Omega$ and R$_1$ = 1\,k$\Omega$, and the divided voltage is fed back to the inverting input. Assume an ideal op-amp operating in its linear region unless stated otherwise. Determine the voltage at the inverting input, the closed-loop gain and the output voltage, and the current in R$_f$. Then simulate.

\labfig{lab4_ex01.png}{9cm}{Exercise 1: non-inverting amplifier}{fig:lab4_ex01}

\theory{8}

\simulation{
\item Draw the circuit of Figure~\ref{fig:lab4_ex01} in LTspice and save it as \file{Lab4/Ex01.asc}.
\item Use the op-amp symbol \cmd{UniversalOpAmp2} from the \textit{[Opamps]} folder (\texttt{F2}). Supply it with two DC sources, $+15$\,V on the V+ pin and $-15$\,V on the V$-$ pin, both referenced to ground.
\item Add the \cmd{.op} analysis (\textit{Simulate $\rightarrow$ Edit Simulation Cmd $\rightarrow$ DC op pnt}), run the simulation and read the node voltages and branch currents.
\item Be careful with signs: LTspice reports the current of a voltage source as positive when it flows \emph{into} its positive terminal, and the current of a two-terminal device as positive when it flows from its first pin to its second pin.
}

\comparetable{
\cmprow{$V_{-}$ (V)}
\cmprow{Gain $V_{out}/V_{in}$}
\cmprow{$V_{out}$ (V)}
\cmprow{$I_{R_f}$ (mA)}
}

\commentlines{2}

\subsection{Exercise 2: Inverting amplifier}
In the circuit of Figure~\ref{fig:lab4_ex02} the 2\,V source is connected to the inverting input through R$_1$ = 1\,k$\Omega$ and the feedback resistor is R$_f$ = 2\,k$\Omega$; the non-inverting input is grounded. Assume an ideal op-amp operating in its linear region unless stated otherwise. Determine the voltage at the inverting input (virtual ground), the gain, the output voltage and the input current. Then simulate.

\labfig{lab4_ex02.png}{9cm}{Exercise 2: inverting amplifier}{fig:lab4_ex02}

\theory{8}

\simulation{
\item Draw the circuit of Figure~\ref{fig:lab4_ex02} in LTspice and save it as \file{Lab4/Ex02.asc}.
\item Use the op-amp symbol \cmd{UniversalOpAmp2} from the \textit{[Opamps]} folder (\texttt{F2}). Supply it with two DC sources, $+15$\,V on the V+ pin and $-15$\,V on the V$-$ pin, both referenced to ground.
\item Add the \cmd{.op} analysis (\textit{Simulate $\rightarrow$ Edit Simulation Cmd $\rightarrow$ DC op pnt}), run the simulation and read the node voltages and branch currents.
\item Be careful with signs: LTspice reports the current of a voltage source as positive when it flows \emph{into} its positive terminal, and the current of a two-terminal device as positive when it flows from its first pin to its second pin.
}

\comparetable{
\cmprow{$V_{-}$ (V)}
\cmprow{Gain $V_{out}/V_{in}$}
\cmprow{$V_{out}$ (V)}
\cmprow{Input current (mA)}
}

\commentlines{2}

\subsection{Exercise 3: Voltage follower}
In the voltage follower of Figure~\ref{fig:lab4_ex03} a 3\,V source is connected to the non-inverting input; the output is connected directly to the inverting input and loaded by R$_L$ = 1\,k$\Omega$. Determine $V_{out}$ and the load current, and explain why the circuit is used as a buffer. Then simulate.

\labfig{lab4_ex03.png}{9cm}{Exercise 3: voltage follower}{fig:lab4_ex03}

\theory{8}

\simulation{
\item Draw the circuit of Figure~\ref{fig:lab4_ex03} in LTspice and save it as \file{Lab4/Ex03.asc}.
\item Use the op-amp symbol \cmd{UniversalOpAmp2} from the \textit{[Opamps]} folder (\texttt{F2}). Supply it with two DC sources, $+15$\,V on the V+ pin and $-15$\,V on the V$-$ pin, both referenced to ground.
\item Add the \cmd{.op} analysis (\textit{Simulate $\rightarrow$ Edit Simulation Cmd $\rightarrow$ DC op pnt}), run the simulation and read the node voltages and branch currents.
\item Add a second load of 100\,$\Omega$ in parallel and check that the output voltage remains the same; read the current delivered by the op-amp in both cases.
}

\comparetable{
\cmprow{$V_{out}$ (V)}
\cmprow{Load current, 1\,k$\Omega$ (mA)}
\cmprow{Load current, 1\,k$\Omega\parallel100\,\Omega$ (mA)}
}

\commentlines{2}

\subsection{Exercise 4: High-current voltage follower}
An op-amp cannot deliver a large current. In the circuit of Figure~\ref{fig:lab4_ex04} an NPN transistor (collector connected to a 12\,V supply) is added at the output of the op-amp as a current booster, and a feedback resistor R$_f$ = 10\,k$\Omega$ connects the emitter to the inverting input; the 3\,V source drives the non-inverting input and R$_L$ = 1\,k$\Omega$ is the load. Determine $V_{out}$, the load current, the base current of the transistor and the output voltage of the op-amp. Assume an ideal op-amp operating in its linear region unless stated otherwise. Then simulate.

\labfig{lab4_ex04.png}{10cm}{Exercise 4: high-current voltage follower}{fig:lab4_ex04}

\theory{10}

\simulation{
\item Draw the circuit of Figure~\ref{fig:lab4_ex04} in LTspice and save it as \file{Lab4/Ex04.asc}.
\item Use the op-amp symbol \cmd{UniversalOpAmp2} from the \textit{[Opamps]} folder (\texttt{F2}). Supply it with two DC sources, $+15$\,V on the V+ pin and $-15$\,V on the V$-$ pin, both referenced to ground.
\item Declare the transistor model \cmd{.model QN NPN(Is=1e-14 Bf=100 Br=1)} (\texttt{S} key) and use \cmd{npn} for Q1.
\item Add the \cmd{.op} analysis (\textit{Simulate $\rightarrow$ Edit Simulation Cmd $\rightarrow$ DC op pnt}), run the simulation and read the node voltages and branch currents.
\item Be careful with signs: LTspice reports the current of a voltage source as positive when it flows \emph{into} its positive terminal, and the current of a two-terminal device as positive when it flows from its first pin to its second pin.
}

\comparetable{
\cmprow{$V_{out}$ (V)}
\cmprow{$I_L$ (mA)}
\cmprow{$I_B$ of Q1 ($\micro$A)}
\cmprow{Op-amp output voltage (V)}
}

\commentlines{2}

\subsection{Exercise 5: Voltage follower with gain}
The circuit of Figure~\ref{fig:lab4_ex05} is a non-inverting amplifier with a current-boosting NPN stage in the loop (V$_{IN}$ = 2\,V, R$_1$ = 1\,k$\Omega$ from the inverting input to ground, R$_2$ = 2\,k$\Omega$ from the emitter to the inverting input, R$_L$ = 100\,$\Omega$, collector supply 12\,V). Determine the gain $1+R_2/R_1$, $V_{out}$ and the load current. Assume an ideal op-amp operating in its linear region unless stated otherwise. Then simulate.

\labfig{lab4_ex05.png}{10cm}{Exercise 5: voltage follower with gain}{fig:lab4_ex05}

\theory{8}

\simulation{
\item Draw the circuit of Figure~\ref{fig:lab4_ex05} in LTspice and save it as \file{Lab4/Ex05.asc}.
\item Use the op-amp symbol \cmd{UniversalOpAmp2} from the \textit{[Opamps]} folder (\texttt{F2}). Supply it with two DC sources, $+15$\,V on the V+ pin and $-15$\,V on the V$-$ pin, both referenced to ground.
\item Declare \cmd{.model QN NPN(Is=1e-14 Bf=100 Br=1)} and use \cmd{npn} for Q1.
\item Add the \cmd{.op} analysis (\textit{Simulate $\rightarrow$ Edit Simulation Cmd $\rightarrow$ DC op pnt}), run the simulation and read the node voltages and branch currents.
\item Be careful with signs: LTspice reports the current of a voltage source as positive when it flows \emph{into} its positive terminal, and the current of a two-terminal device as positive when it flows from its first pin to its second pin.
}

\comparetable{
\cmprow{Gain}
\cmprow{$V_{out}$ (V)}
\cmprow{$I_L$ (mA)}
}

\commentlines{2}

\subsection{Exercise 6: Summing amplifiers}
(a) In the inverting summing amplifier of Figure~\ref{fig:lab4_ex06} the three inputs V1 = 1\,V, V2 = 2\,V and V3 = 3\,V are connected through three 10\,k$\Omega$ resistors to the inverting input, and R$_f$ = 10\,k$\Omega$. Determine $V_{out}$. (b) In the non-inverting summing amplifier (Figure~\ref{fig:lab4_ex06b}) the same inputs are connected through 10\,k$\Omega$ resistors to the non-inverting input, with R$_i$ = 10\,k$\Omega$ and R$_f$ = 20\,k$\Omega$. Determine $V_{out}$. Assume an ideal op-amp operating in its linear region unless stated otherwise. Simulate both circuits.

\labfig{lab4_ex06.png}{9cm}{Exercise 6 (a): inverting summing amplifier}{fig:lab4_ex06}

\labfig{lab4_ex06b.png}{9cm}{Exercise 6 (b): non-inverting summing amplifier}{fig:lab4_ex06b}

\theory{10}

\simulation{
\item Draw the circuit of Figure~\ref{fig:lab4_ex06} in LTspice and save it as \file{Lab4/Ex06.asc}.
\item Use the op-amp symbol \cmd{UniversalOpAmp2} from the \textit{[Opamps]} folder (\texttt{F2}). Supply it with two DC sources, $+15$\,V on the V+ pin and $-15$\,V on the V$-$ pin, both referenced to ground.
\item Add the \cmd{.op} analysis (\textit{Simulate $\rightarrow$ Edit Simulation Cmd $\rightarrow$ DC op pnt}), run the simulation and read the node voltages and branch currents.
\item Be careful with signs: LTspice reports the current of a voltage source as positive when it flows \emph{into} its positive terminal, and the current of a two-terminal device as positive when it flows from its first pin to its second pin.
}

\comparetable{
\cmprow{(a) $V_{out}$ (V)}
\cmprow{(b) $V_{in}$ at the non-inverting input (V)}
\cmprow{(b) $V_{out}$ (V)}
}

\commentlines{2}

\subsection{Exercise 7: Low-pass filter}
The circuit of Figure~\ref{fig:lab4_ex07} is a first-order active low-pass filter: R$_1$ = 1\,k$\Omega$ at the input, and the feedback network is R$_f$ = 10\,k$\Omega$ in parallel with C$_f$ = 15.9\,nF (inverting configuration). Determine the DC gain, the cut-off frequency $f_c=1/(2\pi R_fC_f)$ and the gain at $f_c$ and at $10f_c$. Then run an AC analysis from 10\,Hz to 100\,kHz and verify.

\labfig{lab4_ex07.png}{10cm}{Exercise 7: active low-pass filter}{fig:lab4_ex07}

\theory{10}

\simulation{
\item Draw the circuit of Figure~\ref{fig:lab4_ex07} in LTspice and save it as \file{Lab4/Ex07.asc}.
\item Use the op-amp symbol \cmd{UniversalOpAmp2} from the \textit{[Opamps]} folder (\texttt{F2}). Supply it with two DC sources, $+15$\,V on the V+ pin and $-15$\,V on the V$-$ pin, both referenced to ground.
\item Set the AC amplitude of the input source to 1\,V and run \cmd{.ac dec 20 10 100k}.
\item Plot \cmd{V(out)} (magnitude in dB and phase) and read the DC gain and the $-3$\,dB frequency with the cursors.
}

\comparetable{
\cmprow{DC gain (dB)}
\cmprow{Cut-off frequency $f_c$ (Hz)}
\cmprow{Gain at $10f_c$ (dB)}
}

\commentlines{2}

\subsection{Exercise 8: High-pass filter}
The circuit of Figure~\ref{fig:lab4_ex08} is a first-order active high-pass filter: an input capacitor C$_1$ = 100\,nF in series with R$_1$ = 1.6\,k$\Omega$ and a feedback resistor R$_f$ = 16\,k$\Omega$ (inverting configuration). Determine the high-frequency gain, the cut-off frequency $f_c=1/(2\pi R_1C_1)$ and the gain at $f_c/10$. Then run an AC analysis from 10\,Hz to 100\,kHz and verify.

\labfig{lab4_ex08.png}{10cm}{Exercise 8: active high-pass filter}{fig:lab4_ex08}

\theory{10}

\simulation{
\item Draw the circuit of Figure~\ref{fig:lab4_ex08} in LTspice and save it as \file{Lab4/Ex08.asc}.
\item Use the op-amp symbol \cmd{UniversalOpAmp2} from the \textit{[Opamps]} folder (\texttt{F2}). Supply it with two DC sources, $+15$\,V on the V+ pin and $-15$\,V on the V$-$ pin, both referenced to ground.
\item Set the AC amplitude of the input source to 1\,V and run \cmd{.ac dec 20 10 100k}; plot \cmd{V(out)} and read the high-frequency gain and the $-3$\,dB frequency.
}

\comparetable{
\cmprow{High-frequency gain (dB)}
\cmprow{Cut-off frequency $f_c$ (Hz)}
\cmprow{Gain at $f_c/10$ (dB)}
}

\commentlines{2}

\subsection{Exercise 9: Schmitt trigger}
The circuit of Figure~\ref{fig:lab4_ex09} is an inverting Schmitt trigger: the input signal is applied to the inverting input and a part of the output is fed back to the non-inverting input through the divider R$_1$ = R$_2$ = 10\,k$\Omega$. With a saturation voltage of about $\pm13.5$\,V (supplies $\pm15$\,V), determine the upper and the lower threshold voltages and the hysteresis width. Then apply a 10\,V peak, 100\,Hz sine wave, run a transient simulation of 40\,ms and measure the thresholds.

\labfig{lab4_ex09.png}{9cm}{Exercise 9: inverting Schmitt trigger}{fig:lab4_ex09}

\theory{10}

\simulation{
\item Draw the circuit of Figure~\ref{fig:lab4_ex09} in LTspice and save it as \file{Lab4/Ex09.asc}.
\item Use the op-amp symbol \cmd{UniversalOpAmp2} from the \textit{[Opamps]} folder (\texttt{F2}). Supply it with two DC sources, $+15$\,V on the V+ pin and $-15$\,V on the V$-$ pin, both referenced to ground.
\item Set the input to SINE (amplitude 10\,V, 100\,Hz) and run \cmd{.tran 0 40m 0 10u}; plot the input and the output on the same axes.
\item With the cursors, read the input voltage at the instants when the output switches.
}

\comparetable{
\cmprow{Upper threshold (V)}
\cmprow{Lower threshold (V)}
\cmprow{Hysteresis (V)}
\cmprow{Output high level (V)}
}

\commentlines{2}

\subsection{Exercise 10: Difference amplifier}
In the difference amplifier of Figure~\ref{fig:lab4_ex10} the inputs are V1 = 1\,V and V2 = 1.5\,V, R$_1$ = R$_3$ = 10\,k$\Omega$ and R$_f$ = R$_4$ = 47\,k$\Omega$ (V1 to the inverting input, V2 to the non-inverting input). Assume an ideal op-amp operating in its linear region unless stated otherwise. Determine the voltage at the non-inverting input, the gain of the differential signal and $V_{out}$. Then simulate and discuss what happens when V1 = V2.

\labfig{lab4_ex10.png}{10cm}{Exercise 10: difference amplifier}{fig:lab4_ex10}

\theory{10}

\simulation{
\item Draw the circuit of Figure~\ref{fig:lab4_ex10} in LTspice and save it as \file{Lab4/Ex10.asc}.
\item Use the op-amp symbol \cmd{UniversalOpAmp2} from the \textit{[Opamps]} folder (\texttt{F2}). Supply it with two DC sources, $+15$\,V on the V+ pin and $-15$\,V on the V$-$ pin, both referenced to ground.
\item Add the \cmd{.op} analysis (\textit{Simulate $\rightarrow$ Edit Simulation Cmd $\rightarrow$ DC op pnt}), run the simulation and read the node voltages and branch currents.
\item Set V1 = V2 = 1.5\,V and run again to see the common-mode behaviour.
}

\comparetable{
\cmprow{$V_{+}$ (V)}
\cmprow{Differential gain}
\cmprow{$V_{out}$ (V)}
\cmprow{$V_{out}$ for V1 = V2 (V)}
}

\commentlines{2}

\newpage
\section{Advanced exercises}

\subsection{Advanced Exercise 1: Gain-bandwidth limit of a real op-amp}
A real op-amp has a limited gain-bandwidth product. In the non-inverting amplifier of Figure~\ref{fig:lab4_adv01} (R$_1$ = 1\,k$\Omega$, R$_f$ = 9\,k$\Omega$ or 99\,k$\Omega$) measure the bandwidth for a closed-loop gain of 10 and of 100.

\labfig{lab4_adv01.png}{9cm}{Advanced Exercise 1: non-inverting amplifier with a real op-amp}{fig:lab4_adv01}

\subsubsection{Theory calculation}\vspace{-3pt}\notesblock For an op-amp with GBW = 1\,MHz, calculate the closed-loop gains and the expected bandwidths $f_{-3dB}\approx GBW/A_{cl}$.\dotlines{4}

\simulation{
\item Use the op-amp \cmd{UniversalOpAmp2} and set its \texttt{GBW} parameter to 1\,Meg in the value field (read the other default parameters in the LTspice help).
\item Define \cmd{.param Rf=9k}, use \cmd{\{Rf\}} for the feedback resistor and add \cmd{.step param Rf list 9k 99k}.
\item Run \cmd{.ac dec 50 100 10Meg} and measure the $-3$\,dB frequency of each curve with the cursors.
}
\begin{center}\renewcommand{\arraystretch}{1.8}\begin{tabular}{|c|c|c|c|}\hline \textbf{Closed-loop gain} & \textbf{$f_{-3dB}$ theory (kHz)} & \textbf{$f_{-3dB}$ LTspice (kHz)} & \textbf{Error (\%)}\\ \hline 10 & & & \\ \hline 100 & & & \\ \hline \end{tabular}\end{center}
\commentlines{2}
\newpage

\subsection{Advanced Exercise 2: Second-order Sallen--Key low-pass filter}
Figure~\ref{fig:lab4_adv02} shows a unity-gain Sallen--Key low-pass filter (R$_1$ = R$_2$ = 10\,k$\Omega$, C$_1$ = 22\,nF from the node between the resistors to the output, C$_2$ = 10\,nF to ground). Determine its cut-off frequency and quality factor and compare its roll-off with that of the first-order filter of Exercise~7.

\labfig{lab4_adv02.png}{10cm}{Advanced Exercise 2: Sallen--Key low-pass filter}{fig:lab4_adv02}

\subsubsection{Theory calculation}\vspace{-3pt}\notesblock Use $f_0=1/(2\pi R\sqrt{C_1C_2})$ and $Q=\tfrac12\sqrt{C_1/C_2}$ (for $R_1=R_2=R$).\dotlines{4}

\simulation{
\item Set the AC amplitude of the source to 1\,V and run \cmd{.ac dec 50 10 100k}.
\item Read $f_{-3dB}$ and the attenuation one decade above $f_0$ with the cursors, and compare with the first-order filter.
}
\begin{center}\renewcommand{\arraystretch}{1.8}\begin{tabular}{|c|c|c|c|}\hline \textbf{Quantity} & \textbf{Theory} & \textbf{LTspice} & \textbf{Error (\%)}\\ \hline $f_0$ (Hz) & & & \\ \hline $Q$ & & & \\ \hline Attenuation at $10f_0$ (dB) & & & \\ \hline \end{tabular}\end{center}
\commentlines{2}
\newpage

\chap{Field Effect Transistor}

\section{Introduction}
In this laboratory you study field-effect transistors: the characteristics of the JFET and of the depletion and enhancement MOSFETs, their DC bias circuits, and two applications of the MOSFET as a switch (a lamp and a DC motor).

\textbf{After this laboratory you will be able to:}
\begin{itemize}
  \item plot the output and transfer characteristics of a JFET and a MOSFET with DC sweeps and nested sweeps;
  \item calculate and verify the bias of self-bias and voltage-divider circuits;
  \item identify the triode and saturation regions from the simulated voltages;
  \item use a MOSFET as a switch and explain the role of the flyback diode;
  \item measure averages and ripples of transient waveforms with \cmd{.meas}.
\end{itemize}

\textbf{Exercises.} Save every schematic as \file{Lab5/ExNN.asc} (advanced exercises: \file{Lab5/AdvNN.asc}). Calculation spaces are blank; fill in the comparison tables with your own theoretical and simulated values.

\section{LTspice guide}

\subsection{JFET and MOSFET symbols}
Press \texttt{F2} and search for \cmd{njf} (n-channel JFET), \cmd{pjf}, \cmd{nmos} (n-channel MOSFET) or \cmd{pmos}. The MOSFET symbol has four terminals (drain, gate, source and bulk): connect the bulk to the source. Declare the model with a SPICE directive and type its name in the value field of the transistor; for a MOSFET the instance parameters \cmd{L} and \cmd{W} are entered in the \textit{SpiceLine} (right-click the symbol), and they are set to equal values here so that $K_p$ gives the transconductance directly.

\labfig{lab5_guide01.png}{10cm}{Placing a JFET and declaring its model}{fig:lab5_guide01}

\subsection{Nested DC sweeps for FET characteristics}
The family of output characteristics is plotted with a nested DC sweep of the drain-source voltage and of the gate-source voltage (see the LTspice guide of Lab~3). Plot the drain current by clicking on the transistor symbol in the schematic.

\labfig{lab5_guide02.png}{12cm}{Nested DC sweep for a JFET}{fig:lab5_guide02}

\subsection{Pulse sources and PWM}
A PWM gate drive is a pulse source: in the voltage source dialog click \textit{Advanced} and choose PULSE, with the low value, the high value, the delay, the rise and fall times, the on time and the period. The average of a signal is measured with \cmd{.meas TRAN name AVG expression FROM=t1 TO=t2}.

\section{Exercises}

\subsection{Exercise 1: JFET output characteristic}
In the circuit of Figure~\ref{fig:lab5_ex01} the gate of an n-channel JFET is connected to the source ($V_{GS}=0$) and the drain voltage $V_{DS}$ is swept. The JFET has $I_{DSS}=10$\,mA and $V_P=-2$\,V ($I_D=I_{DSS}(1-V_{GS}/V_P)^2$ in saturation). Sketch the expected curve $I_D=f(V_{DS})$, indicate the ohmic region and the saturation region and the value of $V_{DS}$ at the boundary. Then plot the curve with a DC sweep of $V_{DS}$ from 0 to 10\,V and read $I_{DSS}$.

\labfig{lab5_ex01.png}{8cm}{Exercise 1: JFET with $V_{GS}=0$}{fig:lab5_ex01}

\theory{6}

\simulation{
\item Draw the circuit of Figure~\ref{fig:lab5_ex01} in LTspice and save it as \file{Lab5/Ex01.asc}.
\item Declare the JFET model with a directive (\texttt{S} key): \cmd{.model JN NJF(Beta=2.5m Vto=-2 Lambda=0.01)} ($I_{DSS}=\beta V_P^2=10$\,mA, $V_P=-2$\,V). Place the transistor (\cmd{njf}) and type \texttt{JN} in its value field.
\item Run \cmd{.dc VDS 0 10 0.05} and plot \cmd{Id(J1)} (click the JFET).
\item Read the saturation current and the knee voltage with the cursors.
}

\comparetable{
\cmprow{$I_{DSS}$ (mA)}
\cmprow{Boundary $V_{DS}$ (V)}
\cmprow{Slope in the ohmic region (mS)}
}

\commentlines{2}

\subsection{Exercise 2: JFET family of curves (nested DC sweep)}
Use the circuit of Figure~\ref{fig:lab5_ex02} to plot the family of output characteristics of the JFET for $V_{GS}=-2$, $-1.5$, $-1$, $-0.5$ and 0\,V with a nested DC sweep. The JFET has $I_{DSS}=10$\,mA and $V_P=-2$\,V ($I_D=I_{DSS}(1-V_{GS}/V_P)^2$ in saturation). Calculate the saturation current for each $V_{GS}$ and compare it with the curves.

\labfig{lab5_ex02.png}{8cm}{Exercise 2: JFET with variable $V_{GS}$ and $V_{DS}$}{fig:lab5_ex02}

\theory{8}

\simulation{
\item Draw the circuit of Figure~\ref{fig:lab5_ex02} in LTspice and save it as \file{Lab5/Ex02.asc}.
\item Declare the JFET model with a directive (\texttt{S} key): \cmd{.model JN NJF(Beta=2.5m Vto=-2 Lambda=0.01)} ($I_{DSS}=\beta V_P^2=10$\,mA, $V_P=-2$\,V). Place the transistor (\cmd{njf}) and type \texttt{JN} in its value field.
\item Use the nested sweep \cmd{.dc VDS 0 10 0.05 VGS -2 0 0.5} (first source $V_{DS}$, second source $V_{GS}$) and plot \cmd{Id(J1)}.
\item Read the saturation current of each curve at $V_{DS}=8$\,V.
}

\comparetable{
\cmprow{$I_D$ at $V_{GS}=0$ (mA)}
\cmprow{$I_D$ at $V_{GS}=-0.5$\,V (mA)}
\cmprow{$I_D$ at $V_{GS}=-1$\,V (mA)}
\cmprow{$I_D$ at $V_{GS}=-1.5$\,V (mA)}
\cmprow{$I_D$ at $V_{GS}=-2$\,V (mA)}
}

\commentlines{2}

\subsection{Exercise 3: JFET transfer characteristic}
In the circuit of Figure~\ref{fig:lab5_ex03} the drain voltage is fixed at 8\,V and the gate-source voltage is swept. The JFET has $I_{DSS}=10$\,mA and $V_P=-2$\,V ($I_D=I_{DSS}(1-V_{GS}/V_P)^2$ in saturation). Calculate the transconductance $g_m=\partial I_D/\partial V_{GS}$ at $V_{GS}=-1$\,V and at $V_{GS}=-0.5$\,V, then plot the transfer characteristic $I_D=f(V_{GS})$ and the pinch-off voltage.

\labfig{lab5_ex03.png}{8cm}{Exercise 3: JFET transfer characteristic}{fig:lab5_ex03}

\theory{8}

\simulation{
\item Draw the circuit of Figure~\ref{fig:lab5_ex03} in LTspice and save it as \file{Lab5/Ex03.asc}.
\item Declare the JFET model with a directive (\texttt{S} key): \cmd{.model JN NJF(Beta=2.5m Vto=-2 Lambda=0.01)} ($I_{DSS}=\beta V_P^2=10$\,mA, $V_P=-2$\,V). Place the transistor (\cmd{njf}) and type \texttt{JN} in its value field.
\item Run \cmd{.dc VGS -3 0 0.01} and plot \cmd{Id(J1)}.
\item Plot the derivative \cmd{d(Id(J1))} (right-click the plot window $\rightarrow$ \textit{Add Trace}, expression \cmd{d(Id(J1))}) or estimate $g_m$ from two points.
}

\comparetable{
\cmprow{$V_P$ (V)}
\cmprow{$I_D$ at $V_{GS}=-1$\,V (mA)}
\cmprow{$g_m$ at $V_{GS}=-1$\,V (mS)}
\cmprow{$g_m$ at $V_{GS}=-0.5$\,V (mS)}
}

\commentlines{2}

\subsection{Exercise 4: JFET self-bias}
In the self-bias circuit of Figure~\ref{fig:lab5_ex04} (V$_{DD}$ = 15\,V, R$_D$ = 1.2\,k$\Omega$, R$_G$ = 1\,M$\Omega$, R$_S$ = 470\,$\Omega$) determine the quiescent values $I_D$, $V_{GS}$, $V_{DS}$ and $V_D$. The JFET has $I_{DSS}=10$\,mA and $V_P=-2$\,V ($I_D=I_{DSS}(1-V_{GS}/V_P)^2$ in saturation). (Write $V_{GS}=-I_DR_S$ and solve the resulting quadratic equation for $I_D$.) Check that the JFET works in saturation. Then simulate.

\labfig{lab5_ex04.png}{8cm}{Exercise 4: self-biased JFET}{fig:lab5_ex04}

\theory{10}

\simulation{
\item Draw the circuit of Figure~\ref{fig:lab5_ex04} in LTspice and save it as \file{Lab5/Ex04.asc}.
\item Declare the JFET model with a directive (\texttt{S} key): \cmd{.model JN NJF(Beta=2.5m Vto=-2 Lambda=0.01)} ($I_{DSS}=\beta V_P^2=10$\,mA, $V_P=-2$\,V). Place the transistor (\cmd{njf}) and type \texttt{JN} in its value field.
\item Add the \cmd{.op} analysis (\textit{Simulate $\rightarrow$ Edit Simulation Cmd $\rightarrow$ DC op pnt}), run the simulation and read the node voltages and branch currents.
\item Be careful with signs: LTspice reports the current of a voltage source as positive when it flows \emph{into} its positive terminal, and the current of a two-terminal device as positive when it flows from its first pin to its second pin.
}

\comparetable{
\cmprow{$I_D$ (mA)}
\cmprow{$V_{GS}$ (V)}
\cmprow{$V_D$ (V)}
\cmprow{$V_{DS}$ (V)}
}

\commentlines{2}

\subsection{Exercise 5: JFET voltage-divider bias}
In the circuit of Figure~\ref{fig:lab5_ex05} (V$_{DD}$ = 15\,V, R1 = 1\,M$\Omega$, R2 = 270\,k$\Omega$, R$_D$ = 2.2\,k$\Omega$, R$_S$ = 1.5\,k$\Omega$) determine the gate voltage $V_G$, and the quiescent values $I_D$, $V_{GS}$, $V_S$, $V_D$ and $V_{DS}$. The JFET has $I_{DSS}=10$\,mA and $V_P=-2$\,V ($I_D=I_{DSS}(1-V_{GS}/V_P)^2$ in saturation). Compare this bias circuit with the self-bias circuit of Exercise~4 and explain why it makes $I_D$ less dependent on the JFET parameters. Verify by simulation.

\labfig{lab5_ex05.png}{8cm}{Exercise 5: JFET with a voltage divider}{fig:lab5_ex05}

\theory{10}

\simulation{
\item Draw the circuit of Figure~\ref{fig:lab5_ex05} in LTspice and save it as \file{Lab5/Ex05.asc}.
\item Declare the JFET model with a directive (\texttt{S} key): \cmd{.model JN NJF(Beta=2.5m Vto=-2 Lambda=0.01)} ($I_{DSS}=\beta V_P^2=10$\,mA, $V_P=-2$\,V). Place the transistor (\cmd{njf}) and type \texttt{JN} in its value field.
\item Add the \cmd{.op} analysis (\textit{Simulate $\rightarrow$ Edit Simulation Cmd $\rightarrow$ DC op pnt}), run the simulation and read the node voltages and branch currents.
\item Be careful with signs: LTspice reports the current of a voltage source as positive when it flows \emph{into} its positive terminal, and the current of a two-terminal device as positive when it flows from its first pin to its second pin.
}

\comparetable{
\cmprow{$V_G$ (V)}
\cmprow{$I_D$ (mA)}
\cmprow{$V_{GS}$ (V)}
\cmprow{$V_D$ (V)}
\cmprow{$V_{DS}$ (V)}
}

\commentlines{2}

\subsection{Exercise 6: Depletion-mode MOSFET}
A depletion-mode n-channel MOSFET has $I_{DSS}=8$\,mA and a threshold voltage $V_T=-3$\,V ($I_D=I_{DSS}(1-V_{GS}/V_T)^2$). In the self-bias circuit of Figure~\ref{fig:lab5_ex06} (V$_{DD}$ = 12\,V, R$_D$ = 1\,k$\Omega$, R$_G$ = 1\,M$\Omega$, R$_S$ = 500\,$\Omega$) determine $I_D$, $V_{GS}$ and $V_{DS}$ and check the saturation condition. Then simulate. Explain why a depletion MOSFET can work with a negative or a positive $V_{GS}$.

\labfig{lab5_ex06.png}{8cm}{Exercise 6: self-biased depletion MOSFET}{fig:lab5_ex06}

\theory{10}

\simulation{
\item Draw the circuit of Figure~\ref{fig:lab5_ex06} in LTspice and save it as \file{Lab5/Ex06.asc}.
\item Declare the model with a negative threshold: \cmd{.model DMOS NMOS(Kp=1.78m Vto=-3 Lambda=0.01)} ($K_p=2I_{DSS}/V_T^2$). Place \cmd{nmos}, connect its bulk to the source and set the instance parameters to $W=L$ (for example \cmd{L=1u W=1u}), so that $K_p$ is used directly.
\item Add the \cmd{.op} analysis (\textit{Simulate $\rightarrow$ Edit Simulation Cmd $\rightarrow$ DC op pnt}), run the simulation and read the node voltages and branch currents.
\item Be careful with signs: LTspice reports the current of a voltage source as positive when it flows \emph{into} its positive terminal, and the current of a two-terminal device as positive when it flows from its first pin to its second pin.
}

\comparetable{
\cmprow{$I_D$ (mA)}
\cmprow{$V_{GS}$ (V)}
\cmprow{$V_D$ (V)}
\cmprow{$V_{DS}$ (V)}
}

\commentlines{2}

\subsection{Exercise 7: Enhancement MOSFET characteristics}
An n-channel enhancement MOSFET has a threshold voltage $V_T=2$\,V and $I_D=K(V_{GS}-V_T)^2$ in saturation with $K=1$\,mA/V$^2$ (and $I_D=2K[(V_{GS}-V_T)V_{DS}-V_{DS}^2/2]$ in the triode region). Using Figure~\ref{fig:lab5_ex07}, plot the output characteristics for $V_{GS}=3$, 4, 5 and 6\,V with a nested DC sweep. Calculate the saturation current for each $V_{GS}$ and the boundary value of $V_{DS}$ between the triode and the saturation region, and compare with the curves.

\labfig{lab5_ex07.png}{8cm}{Exercise 7: enhancement MOSFET characteristics}{fig:lab5_ex07}

\theory{8}

\simulation{
\item Draw the circuit of Figure~\ref{fig:lab5_ex07} in LTspice and save it as \file{Lab5/Ex07.asc}.
\item Declare \cmd{.model NM NMOS(Kp=2m Vto=2 Lambda=0.01)} and set the instance parameters to $W=L$ (for example \cmd{L=1u W=1u}).
\item Run the nested sweep \cmd{.dc VDS 0 10 0.05 VGS 3 6 1} and plot \cmd{Id(M1)}.
\item Read the saturation current of each curve at $V_{DS}=8$\,V.
}

\comparetable{
\cmprow{$I_D$, $V_{GS}=3$\,V (mA)}
\cmprow{$I_D$, $V_{GS}=4$\,V (mA)}
\cmprow{$I_D$, $V_{GS}=5$\,V (mA)}
\cmprow{$I_D$, $V_{GS}=6$\,V (mA)}
}

\commentlines{2}

\subsection{Exercise 8: Enhancement MOSFET bias}
In the voltage-divider bias circuit of Figure~\ref{fig:lab5_ex08} (V$_{DD}$ = 12\,V, R1 = 1.5\,M$\Omega$, R2 = 1\,M$\Omega$, R$_D$ = 2.2\,k$\Omega$, R$_S$ = 1\,k$\Omega$; $V_T=2$\,V, $K=1$\,mA/V$^2$) determine $V_G$, $I_D$, $V_{GS}$, $V_D$ and $V_{DS}$ and check the saturation condition $V_{DS}>V_{GS}-V_T$. Then simulate.

\labfig{lab5_ex08.png}{8cm}{Exercise 8: enhancement MOSFET with a voltage divider}{fig:lab5_ex08}

\theory{10}

\simulation{
\item Draw the circuit of Figure~\ref{fig:lab5_ex08} in LTspice and save it as \file{Lab5/Ex08.asc}.
\item Declare \cmd{.model NM NMOS(Kp=2m Vto=2 Lambda=0.01)} and set \cmd{L=1u W=1u} in the instance parameters.
\item Add the \cmd{.op} analysis (\textit{Simulate $\rightarrow$ Edit Simulation Cmd $\rightarrow$ DC op pnt}), run the simulation and read the node voltages and branch currents.
\item Be careful with signs: LTspice reports the current of a voltage source as positive when it flows \emph{into} its positive terminal, and the current of a two-terminal device as positive when it flows from its first pin to its second pin.
}

\comparetable{
\cmprow{$V_G$ (V)}
\cmprow{$I_D$ (mA)}
\cmprow{$V_{GS}$ (V)}
\cmprow{$V_D$ (V)}
\cmprow{$V_{DS}$ (V)}
}

\commentlines{2}

\subsection{Exercise 9: MOSFET as a switch}
A MOSFET can switch a load such as a lamp (modelled by a 47\,$\Omega$ resistor connected to a 12\,V supply, see Figure~\ref{fig:lab5_ex09}). The gate is driven by a 0\,V/5\,V pulse train (period 20\,ms, 50\,\% duty cycle). The transistor has $V_T=2$\,V and $K_p=0.5$\,A/V$^2$ (on-resistance in the triode region $R_{DS(on)}\approx1/[K_p(V_{GS}-V_T)]$). Calculate $R_{DS(on)}$, the load current, the drain voltage and the power dissipated in the transistor when it is on, and the drain voltage when it is off. Then run a transient simulation of 60\,ms.

\labfig{lab5_ex09.png}{8cm}{Exercise 9: MOSFET switch}{fig:lab5_ex09}

\theory{10}

\simulation{
\item Draw the circuit of Figure~\ref{fig:lab5_ex09} in LTspice and save it as \file{Lab5/Ex09.asc}.
\item Declare \cmd{.model NMSW NMOS(Kp=0.5 Vto=2)} and set \cmd{L=1u W=1u}.
\item Set V2 to a pulse source (PULSE: 0\,V, 5\,V, 0 delay, 1\,$\micro$s rise and fall, 10\,ms on time, 20\,ms period) and run \cmd{.tran 0 60m 0 10u}.
\item Plot the gate voltage and the drain voltage and read the on-state drain voltage and the current in R$_L$.
}

\comparetable{
\cmprow{$R_{DS(on)}$ ($\Omega$)}
\cmprow{Load current when on (mA)}
\cmprow{Drain voltage when on (mV)}
\cmprow{Power in the MOSFET when on (mW)}
\cmprow{Drain voltage when off (V)}
}

\commentlines{2}

\subsection{Exercise 10: DC motor control with a MOSFET}
A small DC motor is modelled in Figure~\ref{fig:lab5_ex10} by a 10\,$\Omega$ resistor in series with a 10\,mH inductor, supplied by 12\,V and switched by a MOSFET driven with a 1\,kHz PWM signal with a 50\,\% duty cycle. A diode D1 (flyback diode) is connected across the motor. Calculate the average voltage applied to the motor, the average motor current, and explain the role of the diode (what would happen to the drain voltage at turn-off without it?). Then run a transient simulation of 30\,ms.

\labfig{lab5_ex10.png}{10cm}{Exercise 10: motor drive with a flyback diode}{fig:lab5_ex10}

\theory{10}

\simulation{
\item Draw the circuit of Figure~\ref{fig:lab5_ex10} in LTspice and save it as \file{Lab5/Ex10.asc}.
\item Reuse the MOSFET model of Exercise~9. Drive the gate with PULSE (0\,V, 5\,V, 0 delay, 1\,$\micro$s rise and fall, 0.5\,ms on time, 1\,ms period). Declare a diode model \cmd{.model D1N4148 D(Is=2.52n Rs=0.568 N=1.752)}.
\item Run \cmd{.tran 0 30m 0 2u}. Plot the motor current \cmd{I(Lm)} and the drain voltage.
\item Measure the average current in steady state with \cmd{.meas TRAN Iavg AVG I(Lm) FROM=20m TO=30m}.
\item Remove the diode and run again: observe the drain voltage spikes.
}

\comparetable{
\cmprow{Average motor voltage (V)}
\cmprow{Average motor current (mA)}
\cmprow{Current ripple (mA)}
\cmprow{Peak drain voltage without diode (V)}
}

\commentlines{2}

\newpage
\section{Advanced exercises}

\subsection{Advanced Exercise 1: On-resistance of a MOSFET versus the gate voltage}
In the circuit of Figure~\ref{fig:lab5_adv01} the drain is held at a small voltage (0.1\,V) so that the MOSFET works in the triode region. Study the on-resistance $R_{DS(on)}=V_{DS}/I_D$ as a function of the gate-source voltage.

\labfig{lab5_adv01.png}{8cm}{Advanced Exercise 1: MOSFET in the triode region}{fig:lab5_adv01}

\subsubsection{Theory calculation}\vspace{-3pt}\notesblock Using $R_{DS(on)}\approx1/[K_p(V_{GS}-V_T)]$ with $K_p=0.5$\,A/V$^2$ and $V_T=2$\,V, calculate $R_{DS(on)}$ for $V_{GS}=3$, 5 and 10\,V.\dotlines{4}

\simulation{
\item Run \cmd{.dc VGS 0 10 0.05} and plot \cmd{0.1/Id(M1)} (right-click the plot $\rightarrow$ \textit{Add Trace}). Limit the vertical axis to see the useful range.
\item Read $R_{DS(on)}$ at $V_{GS}=3$, 5 and 10\,V with the cursors.
}
\begin{center}\renewcommand{\arraystretch}{1.8}\begin{tabular}{|c|c|c|c|}\hline \textbf{$V_{GS}$ (V)} & \textbf{$R_{DS(on)}$ theory ($\Omega$)} & \textbf{$R_{DS(on)}$ LTspice ($\Omega$)} & \textbf{Error (\%)}\\ \hline 3 & & & \\ \hline 5 & & & \\ \hline 10 & & & \\ \hline \end{tabular}\end{center}
\commentlines{2}
\newpage

\subsection{Advanced Exercise 2: CMOS inverter}
Figure~\ref{fig:lab5_adv02} shows a CMOS inverter made of an n-channel MOSFET ($V_T=2$\,V, $K_p=2$\,mA/V$^2$) and a p-channel MOSFET ($V_T=-2$\,V, $K_p=1$\,mA/V$^2$), supplied by 5\,V. Plot its voltage transfer characteristic and determine the logic switching threshold.

\labfig{lab5_adv02.png}{8cm}{Advanced Exercise 2: CMOS inverter}{fig:lab5_adv02}

\subsubsection{Theory calculation}\vspace{-3pt}\notesblock Explain qualitatively which transistor conducts for $V_{in}=0$\,V and $V_{in}=5$\,V and estimate the input voltage at which both transistors are in saturation.\dotlines{4}

\simulation{
\item Declare \cmd{.model NM NMOS(Kp=2m Vto=2 Lambda=0.01)} and \cmd{.model PM PMOS(Kp=1m Vto=-2 Lambda=0.01)}; use \cmd{nmos} and \cmd{pmos} symbols with $W=L$.
\item Run \cmd{.dc Vin 0 5 0.01} and plot \cmd{V(out)}; read the threshold where $V_{out}=V_{in}$.
\item Plot the supply current \cmd{I(V1)} to see the current peak during the transition.
}
\begin{center}\renewcommand{\arraystretch}{1.8}\begin{tabular}{|c|c|c|c|}\hline \textbf{Quantity} & \textbf{Theory} & \textbf{LTspice} & \textbf{Error (\%)}\\ \hline Switching threshold (V) & & & \\ \hline Output for $V_{in}=0$\,V (V) & & & \\ \hline Output for $V_{in}=5$\,V (V) & & & \\ \hline \end{tabular}\end{center}
\commentlines{2}
\newpage

\chap{Review Laboratory: Circuits with Diodes and Transistors}

\section{Introduction}
This laboratory is a review of the first five. It contains typical problems of a mid-term project: equivalent resistance, Y--$\Delta$ transformation, circuits with diodes and Zener diodes, and circuits with bipolar transistors. You can use it to prepare for the mid-term examination.

\textbf{After this laboratory you will be able to:}
\begin{itemize}
  \item apply Y--$\Delta$ transformations and verify them by simulation;
  \item analyse circuits with diodes and Zener diodes by deciding the conduction state;
  \item analyse BJT circuits and detect saturation with a DC sweep;
  \item combine Zener diodes and transistors in a regulator and study its load regulation;
  \item use \cmd{.step}, \cmd{mc()} and \cmd{.tf} to automate the analysis.
\end{itemize}

\textbf{Exercises.} Save every schematic as \file{Lab6/ExNN.asc} (advanced exercises: \file{Lab6/AdvNN.asc}). Calculation spaces are blank; fill in the comparison tables with your own theoretical and simulated values.

\section{LTspice guide}

\subsection{Strategy for a review laboratory}
This laboratory brings together the previous ones: resistor networks, diodes, Zener diodes and bipolar transistors. For each exercise, first decide the operating state of every non-linear device (is the diode conducting? is the Zener in breakdown? is the transistor active or saturated?) and write down your assumptions, then calculate, then simulate. A wrong assumption shows up as an inconsistent result (for example a negative diode current).

\subsection{Useful LTspice commands}
\textbf{DC sweep}: \cmd{.dc V1 0 6 0.05}. \textbf{Parameters and steps}: \cmd{.param RL=100}, \cmd{.step param RL list 50 100 1k}. \textbf{Transient}: \cmd{.tran 0 4m 0 5u}. \textbf{Measurements}: \cmd{.meas TRAN Vmax MAX V(out)}. \textbf{Random values}: \cmd{mc(value,tolerance)} used with \cmd{.step param run 1 20 1}. \textbf{Thevenin equivalent}: \cmd{.tf V(out) V1}.

\subsection{Checking a bias point}
After a \cmd{.op} run, check each non-linear device: a conducting diode has about 0.6--0.8\,V across it; a Zener in breakdown has a voltage close to $V_Z$ and a reverse current; a transistor in the active region has $V_{CE}>0.2$\,V and $I_C\approx\beta I_B$.

\section{Exercises}

\subsection{Exercise 1: Equivalent resistance of a bridge network}
In the circuit of Figure~\ref{fig:lab6_ex01} a 24\,V source feeds a 4\,$\Omega$ resistor in series with a bridge made of 20, 30, 60 and 40\,$\Omega$ resistors and a 10\,$\Omega$ resistor between the middle nodes. Use a Y--$\Delta$ transformation to find the total equivalent resistance seen by the source, the current delivered by the source and the power dissipated in the 10\,$\Omega$ resistor. Then simulate.

\labfig{lab6_ex01.png}{10cm}{Exercise 1: bridge network}{fig:lab6_ex01}

\theory{10}

\simulation{
\item Draw the circuit of Figure~\ref{fig:lab6_ex01} in LTspice and save it as \file{Lab6/Ex01.asc}.
\item Add the \cmd{.op} analysis (\textit{Simulate $\rightarrow$ Edit Simulation Cmd $\rightarrow$ DC op pnt}), run the simulation and read the node voltages and branch currents.
\item $R_{eq}=V_1/I$, where $I$ is the magnitude of the source current.
\item Compute the power in the 10\,$\Omega$ resistor from its voltage and current.
\item Be careful with signs: LTspice reports the current of a voltage source as positive when it flows \emph{into} its positive terminal, and the current of a two-terminal device as positive when it flows from its first pin to its second pin.
}

\comparetable{
\cmprow{$R_{eq}$ ($\Omega$)}
\cmprow{$I$ (A)}
\cmprow{$P_{10\,\Omega}$ (W)}
}

\commentlines{2}

\subsection{Exercise 2: Equivalent resistance with equal resistors}
Five equal resistors R = 1\,k$\Omega$ are connected as in Figure~\ref{fig:lab6_ex02}: two in the upper branch, two in the lower branch and one in the middle between the two intermediate nodes. Find the equivalent resistance between A and B and the current in the middle resistor, and explain the result (what is the condition for a balanced bridge?). A 10\,V source is connected between A and B for the simulation.

\labfig{lab6_ex02.png}{8cm}{Exercise 2: network of equal resistors}{fig:lab6_ex02}

\theory{8}

\simulation{
\item Draw the circuit of Figure~\ref{fig:lab6_ex02} in LTspice and save it as \file{Lab6/Ex02.asc}.
\item Add the \cmd{.op} analysis (\textit{Simulate $\rightarrow$ Edit Simulation Cmd $\rightarrow$ DC op pnt}), run the simulation and read the node voltages and branch currents.
\item Read the current of the middle resistor and the voltage across it.
\item Be careful with signs: LTspice reports the current of a voltage source as positive when it flows \emph{into} its positive terminal, and the current of a two-terminal device as positive when it flows from its first pin to its second pin.
}

\comparetable{
\cmprow{$R_{AB}$ (k$\Omega$)}
\cmprow{Current in the middle resistor (mA)}
\cmprow{Voltage across the middle resistor (V)}
}

\commentlines{2}

\subsection{Exercise 3: Current calculation with a delta network}
In the circuit of Figure~\ref{fig:lab6_ex03} a 6\,mA current source feeds node A. The resistors are R1 = 6\,k$\Omega$ (A--B), R2 = 3\,k$\Omega$ (B--C), R3 = 12\,k$\Omega$ (A--C), R4 = 4\,k$\Omega$ (B to ground) and R5 = 2\,k$\Omega$ (C to ground). Use a $\Delta$--Y transformation (or nodal analysis) to calculate the three node voltages and the current in every resistor. Then simulate.

\labfig{lab6_ex03.png}{9cm}{Exercise 3: delta network with a current source}{fig:lab6_ex03}

\theory{12}

\simulation{
\item Draw the circuit of Figure~\ref{fig:lab6_ex03} in LTspice and save it as \file{Lab6/Ex03.asc}.
\item Add the \cmd{.op} analysis (\textit{Simulate $\rightarrow$ Edit Simulation Cmd $\rightarrow$ DC op pnt}), run the simulation and read the node voltages and branch currents.
\item Be careful with signs: LTspice reports the current of a voltage source as positive when it flows \emph{into} its positive terminal, and the current of a two-terminal device as positive when it flows from its first pin to its second pin.
}

\comparetable{
\cmprow{$V_A$ (V)}
\cmprow{$V_B$ (V)}
\cmprow{$V_C$ (V)}
\cmprow{$I_{R1}$ (mA)}
\cmprow{$I_{R2}$ (mA)}
\cmprow{$I_{R3}$ (mA)}
}

\commentlines{2}

\subsection{Exercise 4: Diode circuit 1}
In the circuit of Figure~\ref{fig:lab6_ex04} a 10\,V source feeds two branches, 2\,k$\Omega$ with 8\,k$\Omega$ and 4\,k$\Omega$ with 6\,k$\Omega$, and a diode is connected between the two middle nodes ($V_1$ at the anode, $V_2$ at the cathode). Assume the diode drops 0.7\,V when it conducts. Determine whether it conducts, then calculate $V_1$, $V_2$ and the currents through all the resistors and through the diode. Then simulate.

\labfig{lab6_ex04.png}{8cm}{Exercise 4: diode in a resistor bridge}{fig:lab6_ex04}

\theory{12}

\simulation{
\item Draw the circuit of Figure~\ref{fig:lab6_ex04} in LTspice and save it as \file{Lab6/Ex04.asc}.
\item Declare the diode model with a directive (\texttt{S} key), for example \cmd{.model D1N4148 D(Is=2.52n Rs=0.568 N=1.752)}, and type the model name in the value field of the diode.
\item Add the \cmd{.op} analysis (\textit{Simulate $\rightarrow$ Edit Simulation Cmd $\rightarrow$ DC op pnt}), run the simulation and read the node voltages and branch currents.
\item Be careful with signs: LTspice reports the current of a voltage source as positive when it flows \emph{into} its positive terminal, and the current of a two-terminal device as positive when it flows from its first pin to its second pin.
}

\comparetable{
\cmprow{$V_1$ (V)}
\cmprow{$V_2$ (V)}
\cmprow{$I_D$ (mA)}
\cmprow{$I_{R3}$ (mA)}
\cmprow{$I_{R4}$ (mA)}
}

\commentlines{2}

\subsection{Exercise 5: Diode circuit 2}
In the circuit of Figure~\ref{fig:lab6_ex05} a 30\,V battery feeds R1 = 1\,k$\Omega$ in series with R2 = 1.5\,k$\Omega$; at the end of R2 the circuit splits into R3 = 4.7\,k$\Omega$ to ground and the series connection of R4 = 4.7\,k$\Omega$ and the diode D1 to ground. Decide whether the diode conducts (use $V_F=0.7$\,V), then determine the current in each resistor and in the diode and the voltage at the node between R2, R3 and R4. Then simulate.

\labfig{lab6_ex05.png}{10cm}{Exercise 5: diode with a resistor network}{fig:lab6_ex05}

\theory{12}

\simulation{
\item Draw the circuit of Figure~\ref{fig:lab6_ex05} in LTspice and save it as \file{Lab6/Ex05.asc}.
\item Declare the diode model with a directive (\texttt{S} key), for example \cmd{.model D1N4148 D(Is=2.52n Rs=0.568 N=1.752)}, and type the model name in the value field of the diode.
\item Add the \cmd{.op} analysis (\textit{Simulate $\rightarrow$ Edit Simulation Cmd $\rightarrow$ DC op pnt}), run the simulation and read the node voltages and branch currents.
\item Be careful with signs: LTspice reports the current of a voltage source as positive when it flows \emph{into} its positive terminal, and the current of a two-terminal device as positive when it flows from its first pin to its second pin.
}

\comparetable{
\cmprow{$I_{R1}$ (mA)}
\cmprow{$I_{R3}$ (mA)}
\cmprow{$I_{D}$ (mA)}
\cmprow{Voltage at the node (V)}
}

\commentlines{2}

\subsection{Exercise 6: Zener diode regulator}
In the circuit of Figure~\ref{fig:lab6_ex06} a 16\,V source feeds, through a series resistor R = 220\,$\Omega$, a Zener diode ($V_Z=12$\,V) in parallel with a load R$_L$. Determine the load voltage and the currents $I$, $I_Z$ and $I_L$ for R$_L$ = 10\,k$\Omega$ and for R$_L$ = 1\,k$\Omega$, and the power dissipated in the Zener diode in each case. What is the minimum load resistance for which the Zener diode still regulates? Then simulate.

\labfig{lab6_ex06.png}{9cm}{Exercise 6: Zener regulator}{fig:lab6_ex06}

\theory{12}

\simulation{
\item Draw the circuit of Figure~\ref{fig:lab6_ex06} in LTspice and save it as \file{Lab6/Ex06.asc}.
\item Declare the Zener model with a directive, for example \cmd{.model DZ12 D(Is=1e-14 Bv=12 Ibv=10m Rs=2)}, and place the diode in the reverse direction (cathode to the positive node).
\item Set the load to \cmd{\{RL\}} with \cmd{.param RL=10k} and run both loads with \cmd{.step param RL list 1k 10k}.
\item Add the \cmd{.op} analysis (\textit{Simulate $\rightarrow$ Edit Simulation Cmd $\rightarrow$ DC op pnt}), run the simulation and read the node voltages and branch currents.
\item Be careful with signs: LTspice reports the current of a voltage source as positive when it flows \emph{into} its positive terminal, and the current of a two-terminal device as positive when it flows from its first pin to its second pin.
}

\comparetable{
\cmprow{$V_L$, 10\,k$\Omega$ (V)}
\cmprow{$I_Z$, 10\,k$\Omega$ (mA)}
\cmprow{$P_Z$, 10\,k$\Omega$ (mW)}
\cmprow{$V_L$, 1\,k$\Omega$ (V)}
\cmprow{$I_Z$, 1\,k$\Omega$ (mA)}
\cmprow{$P_Z$, 1\,k$\Omega$ (mW)}
}

\commentlines{2}

\subsection{Exercise 7: BJT driven into saturation}
In the circuit of Figure~\ref{fig:lab6_ex07} an input source V1 drives the base of an NPN transistor through R1 = 10\,k$\Omega$; the collector is connected to a 10\,V supply through R2 = 470\,$\Omega$ and the emitter is grounded. Assume $\beta=100$, $V_{BE}=0.7$\,V and $V_{CE(sat)}=0.2$\,V. Determine the input voltage $V_1$ at which the transistor enters saturation and the collector current at that point. Then sweep V1 from 0 to 6\,V with a DC sweep and plot $V_C$.

\labfig{lab6_ex07.png}{9cm}{Exercise 7: BJT driven by a swept source}{fig:lab6_ex07}

\theory{10}

\simulation{
\item Draw the circuit of Figure~\ref{fig:lab6_ex07} in LTspice and save it as \file{Lab6/Ex07.asc}.
\item Declare \cmd{.model QN NPN(Is=1e-14 Bf=100 Br=1)} and type \texttt{QN} in the value field of the transistor.
\item Run \cmd{.dc V1 0 6 0.05} and plot \cmd{V(c)} and \cmd{Ic(Q1)}.
\item Find the value of $V_1$ at which the collector voltage reaches the saturation value (cursors).
}

\comparetable{
\cmprow{$V_1$ at saturation (V)}
\cmprow{$I_C$ at saturation (mA)}
\cmprow{$V_C$ for $V_1=1$\,V (V)}
\cmprow{$V_C$ for $V_1=6$\,V (V)}
}

\commentlines{2}

\subsection{Exercise 8: BJT with a Zener reference}
In the circuit of Figure~\ref{fig:lab6_ex08} a Zener diode with $V_Z=1.7$\,V keeps the base of an NPN transistor at a fixed voltage; V$_{CC}$ = 6\,V, R$_B$ = 8.6\,k$\Omega$, R$_L$ = 400\,$\Omega$ in the collector, R$_E$ = 100\,$\Omega$ in the emitter. Assume $\beta=100$ and $V_{BE}=0.7$\,V. Determine $V_E$, $I_E$, $I_C$, $I_B$, $V_{CE}$, the Zener current and the collector-base voltage, and verify that the transistor works in the active region. Then simulate.

\labfig{lab6_ex08.png}{9cm}{Exercise 8: BJT with a Zener diode at the base}{fig:lab6_ex08}

\theory{12}

\simulation{
\item Draw the circuit of Figure~\ref{fig:lab6_ex08} in LTspice and save it as \file{Lab6/Ex08.asc}.
\item Declare \cmd{.model QN NPN(Is=1e-14 Bf=100 Br=1)} and type \texttt{QN} in the value field of the transistor.
\item Declare a low-voltage Zener model, for example \cmd{.model DZ17 D(Is=1e-14 Bv=1.7 Ibv=1m Rs=2)}.
\item Add the \cmd{.op} analysis (\textit{Simulate $\rightarrow$ Edit Simulation Cmd $\rightarrow$ DC op pnt}), run the simulation and read the node voltages and branch currents.
\item Be careful with signs: LTspice reports the current of a voltage source as positive when it flows \emph{into} its positive terminal, and the current of a two-terminal device as positive when it flows from its first pin to its second pin.
}

\comparetable{
\cmprow{$V_E$ (V)}
\cmprow{$I_C$ (mA)}
\cmprow{$I_B$ ($\micro$A)}
\cmprow{$I_Z$ (mA)}
\cmprow{$V_{CE}$ (V)}
}

\commentlines{2}

\subsection{Exercise 9: Diode clipper}
The circuit of Figure~\ref{fig:lab6_ex09} is a diode clipper: a 10\,V peak, 1\,kHz sine source feeds a 1\,k$\Omega$ resistor; the output node is connected through a diode D1 to a 3\,V reference (anode at the output) and through a diode D2 to ground (cathode at the output). With $V_F=0.7$\,V determine the maximum and the minimum output voltage and sketch the output. Then run a transient simulation of 4\,ms.

\labfig{lab6_ex09.png}{9cm}{Exercise 9: diode clipper}{fig:lab6_ex09}

\theory{8}

\simulation{
\item Draw the circuit of Figure~\ref{fig:lab6_ex09} in LTspice and save it as \file{Lab6/Ex09.asc}.
\item Declare the diode model with a directive (\texttt{S} key), for example \cmd{.model D1N4148 D(Is=2.52n Rs=0.568 N=1.752)}, and type the model name in the value field of the diode.
\item Set the source to SINE (amplitude 10\,V, 1\,kHz) and run \cmd{.tran 0 4m 0 5u}; plot the input and the output.
\item Read the maximum and minimum of the output with \cmd{.meas TRAN Vmax MAX V(out)} and \cmd{.meas TRAN Vmin MIN V(out)} or with the cursors.
}

\comparetable{
\cmprow{Maximum output voltage (V)}
\cmprow{Minimum output voltage (V)}
}

\commentlines{2}

\subsection{Exercise 10: Series regulator with a Zener diode and a BJT}
In the series regulator of Figure~\ref{fig:lab6_ex10} a 12\,V source feeds a Zener diode ($V_Z=5.1$\,V) through R1 = 1\,k$\Omega$; the Zener voltage is applied to the base of an NPN transistor (collector connected to 12\,V) whose emitter is the output, loaded by R$_L$. Assume $\beta=100$ and $V_{BE}=0.7$\,V. Determine the output voltage, the Zener current, the load current and the base current for R$_L$ = 100\,$\Omega$. Explain how the regulator keeps the output constant when the load changes from 50\,$\Omega$ to 1\,k$\Omega$ and why the transistor makes it better than a simple Zener regulator. Then simulate.

\labfig{lab6_ex10.png}{9cm}{Exercise 10: series regulator}{fig:lab6_ex10}

\theory{12}

\simulation{
\item Draw the circuit of Figure~\ref{fig:lab6_ex10} in LTspice and save it as \file{Lab6/Ex10.asc}.
\item Declare \cmd{.model QN NPN(Is=1e-14 Bf=100 Br=1)} and type \texttt{QN} in the value field of the transistor.
\item Declare \cmd{.model DZ51 D(Is=1e-14 Bv=5.1 Ibv=10m Rs=2)}.
\item Set the load to \cmd{\{RL\}}, define \cmd{.param RL=100} and run \cmd{.step param RL list 50 100 1k}.
\item Add the \cmd{.op} analysis (\textit{Simulate $\rightarrow$ Edit Simulation Cmd $\rightarrow$ DC op pnt}), run the simulation and read the node voltages and branch currents.
\item Be careful with signs: LTspice reports the current of a voltage source as positive when it flows \emph{into} its positive terminal, and the current of a two-terminal device as positive when it flows from its first pin to its second pin.
}

\comparetable{
\cmprow{$V_{out}$, $R_L=50\,\Omega$ (V)}
\cmprow{$V_{out}$, $R_L=100\,\Omega$ (V)}
\cmprow{$V_{out}$, $R_L=1\,$k$\Omega$ (V)}
\cmprow{$I_Z$, $R_L=100\,\Omega$ (mA)}
\cmprow{$I_B$, $R_L=100\,\Omega$ ($\micro$A)}
}

\commentlines{2}

\newpage
\section{Advanced exercises}

\subsection{Advanced Exercise 1: Monte Carlo tolerance analysis}
The resistors of the regulator of Exercise~10 (Figure~\ref{fig:lab6_adv01}) have a 5\,\% tolerance. Study the spread of the output voltage over 20 random draws of the resistor values.

\labfig{lab6_adv01.png}{9cm}{Advanced Exercise 1: series regulator with tolerance}{fig:lab6_adv01}

\subsubsection{Theory calculation}\vspace{-3pt}\notesblock Which resistor affects the output voltage more strongly, R1 or R$_L$? Estimate the worst-case spread of $V_{out}$.\dotlines{4}

\simulation{
\item Give the resistors random values with the function \cmd{mc}: R1 = \cmd{\{mc(1k,0.05)\}} and R$_L$ = \cmd{\{mc(100,0.05)\}}.
\item Add \cmd{.step param run 1 20 1} to repeat the simulation 20 times and \cmd{.op}.
\item Read the output voltage of the 20 runs in the SPICE Error Log (for example \cmd{.meas OP Vout FIND V(out) AT=0}) and fill in the table.
}
\begin{center}\renewcommand{\arraystretch}{1.8}\begin{tabular}{|c|c|}\hline \textbf{Quantity} & \textbf{Value (V)}\\ \hline Minimum $V_{out}$ & \\ \hline Maximum $V_{out}$ & \\ \hline Average $V_{out}$ & \\ \hline Spread (max $-$ min) & \\ \hline \end{tabular}\end{center}
\commentlines{2}
\newpage

\subsection{Advanced Exercise 2: Thevenin equivalent with the \texttt{.tf} command}
Remove R5 from the bridge network of Exercise~1 and consider the terminals where R5 was connected (Figure~\ref{fig:lab6_adv02}). Find the Thevenin equivalent seen from these terminals and use it to calculate the current in R5 = 40\,$\Omega$.

\labfig{lab6_adv02.png}{9cm}{Advanced Exercise 2: network seen by a load resistor}{fig:lab6_adv02}

\subsubsection{Theory calculation}\vspace{-3pt}\notesblock Determine $V_{th}$ (open-circuit voltage) and $R_{th}$ (with the source replaced by a short circuit), then the current in a 40\,$\Omega$ load.\dotlines{6}

\simulation{
\item Run the directive \cmd{.tf V(r) V1} (output node r, input source V1). The SPICE Error Log lists the gain, the input impedance and the output impedance at the node.
\item $V_{th}=\text{gain}\times V_1$ and $R_{th}$ is the output impedance. Check the result by connecting the 40\,$\Omega$ resistor and running \cmd{.op}.
}
\begin{center}\renewcommand{\arraystretch}{1.8}\begin{tabular}{|c|c|c|c|}\hline \textbf{Quantity} & \textbf{Theory} & \textbf{LTspice} & \textbf{Error (\%)}\\ \hline $V_{th}$ (V) & & & \\ \hline $R_{th}$ ($\Omega$) & & & \\ \hline Load current (mA) & & & \\ \hline \end{tabular}\end{center}
\commentlines{2}
\newpage


\end{document}
